\documentclass{lmcs} %%% last changed 2014-08-20

%% optional lists of keywords
\keywords{MSO-transduction, MSO-definability, graph decomposisions}

%% read in additional TeX-packages or personal macros here:
%% e.g. \usepackage{tikz}


%%%%%% packages
\usepackage{amsmath,amssymb}
\usepackage{hyperref}
\usepackage{enumitem}
\usepackage{xspace}
\usepackage{footmisc}
\usepackage{tikz}
\usetikzlibrary{arrows.meta,shapes}

%%%%%%% commands
\newcommand*\defd[2][\em\color{red!25!black}]{{#1#2}}
\newcommand*\strong[2][\bfseries\mathversion{bold}]{{#1#2}}

\newcommand{\compl}[1]{\ensuremath{\overline{#1}}}
\newcommand{\arity}{\operatorname{ar}}
\newcommand*{\fullsucc}[3][\to]{\ensuremath{#2\text{$#1$}#3}}
\def\nfullsucc{\fullsucc[\not\to]}
\newcommand{\FO}{\ensuremath{\operatorname{FO}}\xspace}
\newcommand{\MSO}{\ensuremath{\operatorname{MSO}}\xspace}
\newcommand{\CMSO}[1][]{\ensuremath{\C[#1]{\!\!}\MSO}\xspace}

\newcommand{\C}[1][]{\ensuremath{\operatorname{C}_{#1}}\xspace}
\newcommand{\even}{\operatorname{even}}
\newcommand{\odd}{\operatorname{odd}}
\newcommand{\depth}{\operatorname{depth}}
\newcommand{\subtree}{\mathsf{subtree}}
\newcommand{\representative}{\mathsf{rep}}
\newcommand{\layer}{\mathsf{layer}}
\newcommand{\leaf}{\mathsf{leaf}}
\newcommand{\copyVertex}{\mathsf{copy}}
\newcommand{\nodeSubtree}{\mathsf{rep}}
\newcommand{\ancestor}{\mathsf{ancestor}}
\newcommand{\betweeness}{\mathsf{betweeness}}
\newcommand{\parentPV}{\mathsf{parentPV}}
\newcommand{\lca}{\mathsf{lca}}
\newcommand{\node}{\ensuremath{\mathsf{node}}\xspace}
\DeclareMathOperator*\symdiff{\bigtriangleup}
\newcommand\struct[1]{\ensuremath{\langle#1\rangle}}

\newcommand{\copyV}[1][]{\ensuremath{\mathsf{copy}_{#1}}\xspace}
\newcommand{\colorV}[1][]{\ensuremath{\mathsf{color}_{#1}}\xspace}
\newcommand{\parent}{\ensuremath{\mathsf{parent}}\xspace}
\newcommand{\edge}{\ensuremath{\mathsf{edge}}\xspace}
\newcommand{\medge}{\ensuremath{\textsf{m}\text-\mathsf{edge}}\xspace}
\newcommand{\tedge}{\ensuremath{\mathsf{t}\text-\mathsf{edge}}\xspace}
\newcommand{\cedge}{\ensuremath{\mathsf{c}\text-\mathsf{edge}}\xspace}
\newcommand{\epsord}{\ensuremath{\varepsilon\text-\mathsf{ord}}\xspace}
\newcommand{\SET}{\ensuremath{\mathsf{SET}}\xspace}
\newcommand{\SSET}{\ensuremath{\mathsf{SET!}}\xspace}
\newcommand{\BIPARTITION}{\ensuremath{\mathsf{BIPART}}\xspace}
\newcommand{\splitSet}{\ensuremath{\mathcal{B}}}
\newcommand{\joinSet}{\ensuremath{\mathcal{B}_\mathsf{join}^G}}
\newcommand{\marker}{\ensuremath{\mathsf{marker}}\xspace}
\newcommand{\minElement}{\ensuremath{\mathsf{MIN}}\xspace}
\newcommand{\maxElement}{\ensuremath{\mathsf{MAX}}\xspace}


\newcommand{\descendant}{\ensuremath{\mathsf{desc}}\xspace}
\newcommand{\child}{\ensuremath{\mathsf{child}}\xspace}
\NewDocumentCommand{\repr}{o d()}{\ensuremath{{\mathsf{repr}\IfValueT{#1}{_{#1}}\IfValueT{#2}{(#2)}}}\xspace}
\newcommand{\isleaf}{\ensuremath{\mathsf{isleaf}}\xspace}
\newcommand{\leafset}{\ensuremath{\mathsf{leafset}}\xspace}

\newcommand\DEGENERATE{\ensuremath{\mathsf{degenerate}}\xspace}
\newcommand\PRIME{\ensuremath{\mathsf{prime}}\xspace}
\newcommand\LINEAR{\ensuremath{\mathsf{linear}}\xspace}
\newcommand\SERIES{\ensuremath{\mathsf{series}}\xspace}
\newcommand\PARALLEL{\ensuremath{\mathsf{parallel}}\xspace}

\newcommand{\inNeighbours}{\operatorname{in}}
\newcommand{\outNeighbours}{\operatorname{out}}


\newcommand{\etc}{\emph{etc\ldots}\xspace}
\newcommand{\resp}{resp.\xspace}
\newcommand{\aka}{aka\xspace}
\newcommand{\card}[1]{\ensuremath{\left|#1\right|}}
\newcommand{\ifof}{if and only if\xspace}
\renewcommand{\implies}{\rightarrow}

\def\eg{{\em e.g.}}
\def\ie{{\em i.e.}}
\def\cf{{\em cf.}}

%% due to the dependence on amsart.cls, \begin{document} has to occur
%% BEFORE the title and author information:

\begin{document}

% If the title is longer than 55 characters, then specify a shorter running title as the optional argument to \title. The running title should be roughyl at most 55 characters:
\title[CMSO-transducing tree-like graph decompositions]{CMSO-transducing tree-like graph decompositions\rsuper*}
\titlecomment{{\lsuper*}An extended abstract of this work was presented at the 42nd International Symposium on Theoretical Aspects of Computer Science (STACS 2025)~\cite{STACSversion}.}
\thanks{Rutger Campbell: Supported by the Institute for Basic Science (IBS-R029-C1).\\
Mamadou Moustapha Kant\'{e}: Supported by the French National Research Agency (ANR-18-CE40-0025-01 and ANR-20-CE48-0002).}	%optional

% affiliations are numbered automatically with a, b, c (see below)
% use the optional argument to indicate the affiliation(s) of each author
% omit the argument if there is only one author, or only one affiliation
\author[R.~Campbell]{Rutger Campbell}[a]
\author[B.~Guillon]{Bruno Guillon}[b]
\author[M.~M.~Kant\'{e}]{Mamadou Moustapha Kant\'{e}\lmcsorcid{https://orcid.org/0000-0003-1838-7744}}[b]
\author[E.~J.~Kim]{Eun Jung Kim\lmcsorcid{https://orcid.org/0000-0002-6824-0516}}[c,d]
\author[N.~K\"{o}hler]{Noleen K\"{o}hler\lmcsorcid{https://orcid.org/0000-0002-1023-6530}}[e]

% affiliation 1 (automatically numbered a)
\address{Discrete Mathematics Group, Institute for Basic Science, Daejeon, Korea}	
\email{rutger@ibs.re.kr}  %optional

% affiliation 2 (automatically numbered b)
\address{Université Clermont Auvergne, Clermont Auvergne INP, LIMOS, CNRS, Clermont-Ferrand, France}	
\email{bruno.guillon@uca.fr, mamadou.kante@uca.fr} 

% affiliation 3 (automatically numbered c)
\address{KAIST, Daejeon, South Korea}	
\email{eunjung.kim@kaist.ac.kr}  

% affiliation 4 (automatically numbered d)
\address{CNRS, France}	  

% affiliation 5 (automatically numbered e)
\address{University of Leeds, Leeds, UK}	
\email{N.Koehler@leeds.ac.uk}  

%% etc.

%% required for running head on odd and even pages, use suitable
%% abbreviations in case of long titles and many authors:

%%%%%%%%%%%%%%%%%%%%%%%%%%%%%%%%%%%%%%%%%%%%%%%%%%%%%%%%%%%%%%%%%%%%%%%%%%%

%% the abstract has to PRECEDE the command \maketitle:
%% be sure not to issue the \maketitle command twice!

\begin{abstract}
  \noindent We show that given a graph $G$ we can \CMSO-transduce its modular decomposition, its split decomposition and its bi-join decomposition. This improves results by Courcelle [Logical Methods in Computer Science, 2006] who gave such transductions using order-invariant \MSO, a strictly more expressive logic than \CMSO. Our methods more generally yield \CMSO[2]-transductions of the canonical decomposition of weakly-partitive set systems and weakly-bipartitive systems of bipartitions.
\end{abstract}

\maketitle


\section{Introduction}
A decomposition of a graph, especially a tree-like decomposition, is a result of recursive separations of a graph
and is extremely useful for investigating combinatorial properties such as colourability, and for algorithm design. Such a decomposition also plays a fundamental role when one wants to understand the relationship between logic and a graph class.
Different notions of the complexity of a separation motivate different ways to decompose, such as tree-decomposition, branch-decomposition, rank-decomposition and carving-decomposition.
Furthermore, some important graph classes can be defined through the tree-like decomposition they admit;
cographs with cotrees and distance-hereditary graphs with split decompositions being prominent examples.


For a logic $\mathsf{L}$,
an $\mathsf{L}$-transduction 
is a non-deterministic map from a class of relational structures to a new class of relational structures using $\mathsf{L}$-formulas. Transducing
a tree-like decomposition is of particular interest.
Notably, transducing a decomposition of a graph implies that any property that is definable using a decomposition, is also definable on graphs that admit such a decomposition.
Moreover, tree-like decompositions can be often represented by labelled trees, for which the equivalence of recognisability by a tree automaton and definability in CMSO is well known~\cite{TMSOLOG1}.
Hence, it is an interesting question to consider what kind of graph decompositions can be transduced using $\mathsf{L}$-transductions for some extension $\mathsf{L}$ of \MSO.

In a series of papers~\cite{TMSOLOG1,TMSOLOF5,Courcelle96,Courcelle06}, Courcelle investigated the relationship between the graph properties that can be defined in an extension of \MSO and the graph properties that can be recognized by a tree automaton.
In particular, for graphs of bounded treewidth, Courcelle's theorem states that any property that is definable in the logic \MSO with modulo counting predicate, denoted \CMSO, can be recognized by a tree automaton \cite{TMSOLOG1}.
Combining this result with the linear time algorithm for computing tree-decompositions \cite{bodlaender1993linear}, yields that \CMSO\ model-checking can be done in linear time on graphs of bounded treewidth.
The converse statement -- whether recognisability by a tree automaton implies definability in \CMSO\ on graphs of bounded treewidth --
was conjectured by Courcelle in \cite{TMSOLOG1} and finally settled by Boja\'nczyk and Pilipczuk \cite{BojanczykP16}. The key step to obtain this result
is obtaining a tree-decomposition of a graph via an \MSO-transduction, a strategy which was initially proposed in \cite{TMSOLOF5} and is now standard.

The obvious next question is whether an analogous result can be proved for graphs of bounded clique-width and for more general combinatorial objects, most notably, matroids representable over a fixed finite field and of bounded branch-width. Due to the above-mentioned strategy, the key challenge is to produce corresponding tree-like decompositions by \MSO-transduction. It is known that clique-width decompositions can be \MSO-transduced for
graphs of bounded linear clique-width \cite{BojanczykGP21}. However, it is unknown if clique-width decompositions can be \MSO-transduced in general. In fact,
this question remained open even for distance-hereditary graphs, which are precisely graphs of rank-width 1 (thus, of constant clique-width). 

Besides tree-decompositions, the problem of transducing cotrees, and in general
hierarchical decompositions such as modular
decompositions and split decompositions were considered in the literature~\cite{Courcelle96,Courcelle99,Courcelle06,Courcelle13}. In \cite{Courcelle96}, Courcelle provides transductions using order-invariant \MSO for cographs and modular decompositions of
graphs of bounded modular width. Order-invariant \MSO allows the use of a linear order of the set of vertices and is more expressive than \CMSO
\cite{GanzowR08}. The applicability of these transductions was later generalized using the framework of weakly-partitive set systems\footnote{Weakly-partitive set systems are set systems enjoying some nice closure properties, which were then used
to show that some set systems allow canonical tree representations, see for instance the thesis by Montgolfier and Rao
\cite{montgolfierThesis,raoThesis}.} to obtain order-invariant
\MSO-transductions of modular and split decompositions \cite{Courcelle06}. It was left as an open question whether one can get rid of the order (see for instance
\cite{CourcelleT12} where an overview of the result on hierarchical decompositions was given).

\subsection{Our results}
In this paper,
we consider decompositions given by separations that do not overlap with any other separations of the same type.
We view separations of a given kind as a `set system'.
A set system consists of a set $U$, the universe, and a set $\mathcal{S}$ of subsets of $U$. Two sets overlap if they have non-empty intersection but neither of them is contained in the other. If no two elements in a set system $(U,\mathcal{S})$ overlap, i.e. the set system is \emph{laminar}, then the subset relation in $(U,\mathcal{S})$ can be described by a rooted tree $T$, called the \emph{laminar tree} of $(U,\mathcal{S})$.
For any set system $(U, \mathcal{S})$ we can look at the subset of strong sets, i.e., sets that do not overlap with any other set, and the laminar tree $T$ they induce. Additionally, we consider systems of bipartitions $\mathcal{B}$ of a universe $U$. Two bipartitions of $U$ overlap if neither side of one of the bipartition is contained in either side of the other bipartition. In case $(U,\mathcal{B})$ has no overlapping bipartitions, then $(U,\mathcal{B})$ can be described by an undirected tree, also called \emph{laminar tree}, in which bipartitions correspond to edge cuts. We can define the concept of strong bipartitions equivalently and consider the laminar tree induced by the strong bipartitions in $(U,\mathcal{B})$.

Given a graph $G$ we can consider the set system $(V(G),\mathcal{M})$ where $\mathcal{M}$ is the set of all modules in $G$, or the system of bipartition $(V(G),\mathcal{S})$ where $\mathcal{S}$ contains all splits in $G$, or the system of bipartitions $(V(G),\mathcal{B})$ where $\mathcal{B}$ contains all bi-joins in $G$. We obtain the modular decomposition/cotree/split decomposition/bi-join decomposition by equipping the laminar tree of the suitable set system/system of bipartitions with some additional structure. The additional structure allows us to recover the graph from the respective decomposition.

Abstractly, the set systems/systems of bipartitions mentioned are instances of \emph{weakly-partitive set systems}/\emph{weakly-bipartitive systems of bipartitions}.
Roughly speaking, if a set system $(U,\mathcal{S})$ is well behaved, i.e. $(U,\mathcal{S})$ is weakly-partitive (definition in \autoref{sec:preliminaries}), then there is a labelling $\lambda$ of $T$ and a partial order $<$ of its nodes such that $(U,\mathcal{S})$ is completely described by $(T,\lambda,<)$ \cite{chein1981partitive,raoThesis}. A similar statement holds for systems of bipartitions \cite{montgolfierThesis}.

We show the following, wherein each item $\lambda$ is a suitable labelling of the nodes of the laminar tree $T$, $<$ is a partial ordering of its nodes, and $F$ is an additional edge relation defined only on pairs of siblings in $T$. A visualization how results depend on each other is given in \autoref{fig:overview}.

\begin{thm}\label{thm:summary}
    There are non-deterministic \CMSO[2]-transductions $\tau_1,\ldots,\tau_7$ such that: 
    \begin{enumerate}
        \item 
           For any laminar set system $(U,\mathcal{S})$, $\tau_1(U, \mathcal{S})$ is non-empty
           and every output in~$\tau_1(U, \mathcal{S})$ is a laminar tree $T$ of $(U,\mathcal{S})$ (\autoref{thm:transduce laminar tree});
        \item 
            For any graph $G$, $\tau_2(G)$ is non-empty
            and every output in~$\tau_2(G)$ is a modular decomposition $(T,F)$ of $G$ (\autoref{thm:transduce modular decomposition});
        \item 
            For any cograph $G$, $\tau_3(G)$ is non-empty
            and every output in~$\tau_3(G)$ is a cotree $(T,\lambda)$ of $G$ (\autoref{thm:transduce cotree});
        \item 
             For any graph $G$, $\tau_4(G)$ is non-empty
             and every output in~$\tau_4(G)$ is a split decomposition $(T,F)$ of $G$ (\autoref{thm:split});
        \item 
             For any graph $G$, $\tau_5(G)$ is non-empty
             and every output in~$\tau_5(G)$ is a bi-join decomposition $(T,F)$ of $G$ (\autoref{thm:bi-join1} and \autoref{thm:bi-join2});
        \item 
            For any weakly-partitive set system $(U,\mathcal{S})$, $\tau_6(U,\mathcal{S})$ is non-empty
            and every output in~$\tau_6(U,\mathcal{S})$ is a weakly-partitive tree $(T,\lambda,<)$ of $(U,\mathcal{S})$ (\autoref{thm:transduce weakly-partitive tree});
        \item 
            For any weakly-bipartitive set system $(U,\mathcal{B})$, $\tau_7(U,\mathcal{B})$ is non-empty
            and every output in~$\tau_7(U,\mathcal{B})$ is a weakly-bipartitive tree $(T,\lambda,<)$ of $(U,\mathcal{B})$ (\autoref{thm:transduce weakly-bipartitive tree}).
    \end{enumerate}
\end{thm}
The key step in obtaining these transductions
is to transduce the laminar tree $T$ of a set system $(U,\mathcal{S})$.
The crux here is to find a suitable representative of each node of $T$ amongst the elements of $U$
and a non-deterministic coloring which allows the assignment of representatives to nodes by means of a \CMSO[2]-formula.
It should be mentioned that a similar result is claimed in the preprint \cite{Boj23},
where a proof sketch designing a \CMSO[3]-transduction is described.
%
Once the laminar tree is obtained,
the additional relations for each decomposition
can be obtained using a deterministic \MSO-transduction.
Notice that for each of these transductions,
there exists an inverse deterministic \MSO-transduction,
namely a transduction that from the tree-like decomposition outputs the original structure. It is worth mentionning that a corollary of
\autoref{thm:summary}(4) is a \CMSO-transduction outputting a rank-decomposition of width $1$ for every graph of rank-width $1$. 


\begin{figure}[omitted: a diagram of this paper's results, not reproduced here]\label{fig:overview}\end{figure}
\bigbreak

\subsection{Organization}
The paper is organized as follows. In \autoref{sec:preliminaries} we introduce terminology and notation needed. In \autoref{sec:transduce laminar tree} we prove \autoref{thm:transduce laminar tree}. In \autoref{sec:transduce modular decomposition} we provide the proof of \autoref{thm:transduce weakly-partitive tree} and \autoref{thm:transduce modular decomposition} and obtain \autoref{thm:transduce cotree}. In \autoref{sec:split} we provide the proofs of \autoref{thm:transduce weakly-bipartitive tree}, \autoref{thm:split}, \autoref{thm:bi-join1} and \autoref{thm:bi-join2}.

\section{Preliminaries}\label{sec:preliminaries}
\subsection{Graphs, trees, set systems}
\subparagraph{Graphs}
We use standard terminology of graph theory,
and we fix some notations.
Given a \emph{directed graph}~$G$,
its sets of \emph{vertices} and \emph{edges}
are denoted by \defd{$V(G)$} and \defd{$E(G)$}, respectively.
We denote by \defd{uv} an edge~$(u,v)\in E(G)$.
An undirected graph is no more than a directed graph
for which~$E(G)$ is symmetric (\ie, $uv\in E(G)\iff vu\in E(G)$).
The notions of \emph{paths}, \emph{connected components}, \etc are defined as usual.
Given a subset~$Z$ of~$V(G)$,
we denote by \defd{$G[Z]$} the sub-graph of~$G$ induced by~$Z$.

\subparagraph{Trees}
A \defd{tree} is a connected undirected graph without cycles.
In the context of trees,
we use a slightly different terminology than for graphs.
In particular, vertices are called~\defd{nodes},
nodes of degree at most~$1$ are called \defd{leaves},
and nodes of degree greater than~$1$ are called \defd{inner}.
The set of leaves is denoted~\defd{$L(T)$};
thus the set of inner nodes is~$V(T)\setminus L(T)$. For a node $t$ of a tree $T$ and a neighbor $s$ of $t$, we denote by $T^t_s$ the connected component of $T-t$ containing $s$.
We sometimes consider \defd{rooted trees},
namely trees with a distinguished node, called the \defd{root}.
Rooted trees enjoy a natural orientation of their edges
toward the root,
which induces the usual notions of
\defd{parent}, \defd{child}, \defd{sibling}, \defd{ancestor} and \defd{descendant}.
Hence, we represent a rooted tree by a set of nodes with an ancestor/descendant relationship (instead of specifying the root).
We use the convention that every node is one of its own ancestors and descendants.
We refer to ancestors (\resp descendants) of a node
that are not the node itself
as \defd{proper ancestors} (\resp \defd{proper descendants}).
For a node~$t$ of a rooted tree~$T$,
we denote by~\defd{$T_t$} the subtree of~$T$ rooted in~$t$
(\ie, the restriction of~$T$ to the set of descendants of~$t$).


\subparagraph{Set systems and laminar trees}
A \defd{set system} is a pair~$(U,\mathcal{S})$
where~$U$ is a finite set, called the \defd{universe},
and~$\mathcal{S}$ is a family of subsets of~$U$
where~$\emptyset\notin\mathcal{S}$,
$U\in\mathcal{S}$,
and~$\{a\}\in\mathcal{S}$ for every~$a\in U$.%
\footnote{%
    Though these restrictions on~$\mathcal{S}$ are not usual for set systems,
    it is convenient for our contribution
    and it does not significantly impact the generality of set systems:
    every family~$\mathcal{F}$ of subsets of~$U$
    can be associated with a set system~$(U,\mathcal{S})$
    where $\mathcal{S}=\big(\mathcal{F}\setminus\{\emptyset\}\big)\cup\big\{U\big\}\cup\big\{ \{a\} \mid a\in U\big\}$.%
}
Two sets~$X$ and~$Y$ \defd{overlap}
if they are neither disjoint nor related by containment.
A set system~$(U,\mathcal{S})$
is said to be \defd{laminar} (\aka \emph{overlap-free})
when no two sets from~$\mathcal{S}$ overlap.
By extension, a set family~$\mathcal{S}$ is \defd{laminar}
if $\left(\bigcup\mathcal{S}, \mathcal{S}\right)$ is a laminar set system
(note this also requires that~$\emptyset\notin\mathcal{S}$, $\bigcup\mathcal{S}\in\mathcal{S}$, and~$\{a\}\in\mathcal{S}$ for every~$a\in\bigcup\mathcal{S}$).

A laminar family~$\mathcal{S}$ of subsets of~$U$
naturally defines a rooted tree
where the nodes are the sets from~$\mathcal{S}$,
the root is~$U$,
and the ancestor relation corresponds to set inclusion.
We call this rooted tree
the \defd{laminar tree of~$U$ induced by $\mathcal{S}$}
(or \defd{laminar tree of~$(U,\mathcal{S})$}).\xspace
In this rooted tree,
the leaves are the singletons~$\{x\}$ for~$x\in U$,
which we identify with the elements themselves.
That is to say, $L(T)=U$.
Laminar trees have the property
that each inner node has at least two children.
Observe also that the size of a laminar tree
is linearly bounded in the size of the universe.

\subsection{Logic and transductions}
We use \emph{relational structures} to model
both graphs and the various tree-like decompositions used in this paper.
In order to concisely model set systems
we use the more general notion of \emph{extended relational structures},
namely \emph{relational structure} extended with set predicate names.
Such structures also naturally arise as outputs of \emph{\MSO-transductions} defined below.
\smallbreak

Define an \defd{(extended) vocabulary} to be a set of symbols,
each being either a \emph{relation} name~$R$ with associated arity~\defd{$\arity(R)\in\mathbb{N}$},
or a \emph{set predicate} name~$P$ with associated arity~\defd{$\arity(P)\in\mathbb{N}$}.
Set predicate names are aimed to describe relations between sets,
\eg, one may have a unary set predicate for selecting finite sets of even size,
or a binary set predicate for selecting pairs of disjoint sets.
We use capital~$R$ or
names starting with a lowercase letter (\eg, $\edge$, $\ancestor$, $\tedge$) for relation names,
and capital~$P$ or uppercase names (\eg, $\SET$, $\C[2]$) for set predicate names.
To refer to an arbitrary symbol of undetermined kind, we use capital~$Q$.
A \defd{relational vocabulary} is an extended vocabulary in which every symbol is a relation name.

Let~$\Sigma$ be a vocabulary.
An \defd{extended relational structure over~$\Sigma$} (\defd{$\Sigma$-structure})
is a structure $\mathbb{A}=\langle U_{\mathbb{A}}, (Q_{\mathbb{A}})_{Q\in \Sigma}\rangle$
consisting on the one hand of a set~\defd{$U_\mathbb{A}$} called \defd{universe},
and on the other hand, for each symbol~$Q$ in~$\Sigma$,
an \defd{interpretation~$Q_{\mathbb{A}}$ of~$Q$},
which is a relation of arity~$\arity(Q)$
either over the universe if~$Q$ is a relation name,
or over the family of subsets of the universe if~$Q$ is a set predicate name.
When~$\Sigma$ is not extended, $\mathbb{A}$ is simply a \defd{relational structure}.

Given a $\Sigma$-structure~$\mathbb{A}$
and,
for some vocabulary~$\Gamma$,
a $\Gamma$-structure~$\mathbb{B}$,
we write~\defd{$\mathbb{A}\sqsubseteq\mathbb{B}$}
if~$\Sigma\subseteq\Gamma$,
$U_{\mathbb{A}}\subseteq U_{\mathbb{B}}$
and for each symbol~$Q$ in~$\Sigma$,
$Q_{\mathbb{A}}=Q_{\mathbb{B}}$.%
\footnote{%
    We require equality here
    (in particular only elements or subsets of~$U_{\mathbb{A}}$ are related in~$Q_{\mathbb{B}}$).
    This differs from classical notions of inclusions of relational structures
    which typically require equality only on the restriction of the universe to~$U_{\mathbb{A}}$,
    \ie, $Q_{\mathbb{A}}=Q_{\mathbb{B}_{/U_{\mathbb{A}}}}$,
    \eg, in order to correspond to induced graphs.%
}
We write~\defd{$\mathbb{A}\sqcup\mathbb{B}$}
to denote the $(\Sigma\cup\Gamma)$-structure
consisting of the universe~$U_{\mathbb{A}}\cup U_{\mathbb{B}}$
and, for each symbol~$Q\in\Sigma\cup\Gamma$,
the interpretation $Q_{\mathbb{A}\sqcup\mathbb{B}}$
which is~$Q_{\mathbb{A}}$, $Q_{\mathbb{B}}$, or~$Q_{\mathbb{A}}\cup Q_{\mathbb{B}}$
according to whether~$Q$ belongs to~$\Sigma\setminus\Gamma$, to~$\Gamma\setminus\Sigma$, or to~$\Sigma\cap\Gamma$.%
\footnote{%
    We do not require~$\mathbb{A}$ and~$\mathbb{B}$ to be disjoint structures
    whence we may have~$\mathbb{A}\not\sqsubseteq\mathbb{A}\sqcup\mathbb{B}$.%
}

To describe properties of (extended) relational structures, we use \defd{monadic second order logic} (\defd{\MSO}) and refer for instance to
\cite{CE09,FunkMN22,Hlineny06,Strozecki11} for the definition of \MSO on extended relational structures such as matroids or set systems in general. This logic allows quantification both over single elements of the universe
and over subsets of the universe. We also use \defd{counting \MSO} (\defd{\CMSO}), which is the extension of~$\MSO$ with, for every positive integer~$p$, a
unary set predicate~$\C[p]$ that checks whether the size of a subset is divisible by~$p$ or not. We only use~$\C[2]$.
%
As usual, lowercase variables indicates first-order-quantified variables,
while uppercase variables indicates monadic-quantified variables.
For a formula~$\phi$, we write, \eg,~$\phi(x,y,X)$ to indicate that the variables~$x$, $y$, and~$X$ belong to the set of \defd{free variables of~$\phi$},
namely, the set of variables occurring in~$\phi$ that are not bound to a quantifier within~$\phi$.
A \defd{sentence} is a formula without free variables.
\smallbreak

We now fix some (extended) vocabularies that we will use.
\begin{description}[nosep]
    \item[Graphs]
        To model both graphs, unrooted trees, and directed graphs,
        we use the relational vocabulary~$\{\edge\}$
        where~$\edge$ is a relation name of arity~$2$.
        A (directed) graph~$G=(V,E)$ is modeled
        as the $\{\edge\}$-structure $\mathbb{G}$
        with universe~$U_{\mathbb{G}}=V$
        and interpretation~$\edge_{\mathbb{G}}=E$.
        In particular if~$G$ is undirected, then $\edge_{\mathbb{G}}$ is symmetric.
    \item[Rooted trees]
        We use the relational vocabulary~$\{\ancestor\}$ to model rooted trees
        where $\ancestor$ is a relation name of arity~$2$.
        A rooted tree~$T$ is modeled
        as the $\{\ancestor\}$-structure $\mathbb{T}$
        with universe~$U_{\mathbb{T}}=V(T)$
        and the interpretation~$\ancestor_{\mathbb{T}}$
        being the set of pairs~$(u,v)$ such that~$u$ is an ancestor of $v$ in $T$.
        It is routine to define \FO-formulas over this vocabulary
        to express the binary relations parent, child, proper ancestor, (proper) descendant,
        as well as the unary relations leaf and root.%
    \item[Set systems]
        To model set systems, we use the extended vocabulary~$\{\SET\}$
        where $\SET$ is a set predicate name of arity~$1$.
        A set system~$S=(U,\mathcal{S})$ is thus naturally modeled
        as the $\{\SET\}$-structure $\mathbb{S}$
        with universe~$U_{\mathbb{S}}=U$
        and interpretation~$\SET_{\mathbb{S}}=\mathcal{S}$.
\end{description}

\paragraph*{Transductions}
Let~$\Sigma$ and~$\Gamma$ be two extended vocabularies.
A \defd{$\Sigma$-to-$\Gamma$ transduction}
is a set~$\tau$ of pairs formed by
a $\Sigma$-structure, call the \defd{input},
and a $\Gamma$-structure, called the \defd{output}.
We write \defd{$\mathbb{B}\in\tau(\mathbb{A})$} when~$(\mathbb{A},\mathbb{B})\in\tau$.
When for every pair~$(\mathbb{A},\mathbb{B})\in\tau$
we have~$\mathbb{A}\sqsubseteq\mathbb{B}$,
we call~$\tau$ an
\defd{overlay transduction}.
Some transductions
can be defined by means of \MSO- or \CMSO[2]-formulas.
This leads to the notion of \defd{\MSO-} and \defd{\CMSO[2]-transductions}.
Following the presentation of~\cite{BojanczykGP21},
for~$\mathsf{L}$ denoting~\MSO or~\CMSO[2],
define an \defd{$\mathsf{L}$-transduction}
to be a transduction obtained
by composing a finite number of \defd{atomic $\mathsf{L}$-transductions}
of the following kinds.
\begin{description}[nosep]
    \item[Filtering]
        An overlay transduction
        specified by an $\mathsf{L}$-sentence~$\phi$ over the input vocabulary~$\Sigma$,
        which discards the inputs that do not satisfy~$\phi$
        and keeps the other unchanged.
        Hence, it defines a partial function
        (actually, a partial identity)
        from $\Sigma$-structures to $\Sigma$-structures.
    \item[Universe restriction]
        A transduction specified by a $\mathsf{L}$-formula~$\phi$ over the input vocabulary~$\Sigma$,
        with one free first-order variable,
        which restricts the universe to those elements that satisfy~$\phi$.
        The output vocabulary is~$\Sigma$
        and the interpretation of every relation (\resp every predicate) in the output structure
        is defined as the restriction of its interpretation in the input structure
        to those tuples of elements satisfying~$\phi$ (\resp tuples of sets of elements that satisfy~$\phi$).
        This defines a total function from $\Sigma$-structures to $\Sigma$-structures.
    \item[Interpretation]
        A transduction specified by a family~$(\phi_Q)_{Q\in\Gamma}$ over the input vocabulary~$\Sigma$
        where~$\Gamma$ is the output vocabulary and each~$\phi_Q$ has $\arity(Q)$ free variables
        which are first-order if~$Q$ is a relation name and monadic if it is a set predicate name.
        The transduction outputs the~$\Gamma$-structure
        that has the same universe as the input structure
        and in which each relation or predicate~$Q$ is interpreted
        as those set of tuples that satisfy~$\phi_Q$.
        This defines a total function from $\Sigma$-structures to $\Gamma$-structures.
    \item[Copying]
        An overlay transduction parametrized by a positive integer~$k$ that
        adds $k$ copies of each element to the universe.
        The output vocabulary consists in the input vocabulary~$\Sigma$
        extended with~$k$ binary relational symbols~$(\copyV[i])_{i\in[k]}$
        interpreted as pairs of elements~$(x,y)$
        saying that ``$y$ is the $i$-th copy of~$x$''.
        The interpretation of the relations (\resp predicates) of the input structure
        are preserved, on original elements.
        This defines a total function from~$\Sigma$-structures to~$\Gamma$-structures,
        where~$\Gamma=\Sigma\cup\{\copyV[i]\mid i\in[k]\}$.
    \item[Colouring]
        An overlay transduction that adds a new unary relation~$\colorV\notin\Sigma$ to the structure.
        Any possible interpretation yields an output;
        indeed the interpretation is chosen non-deterministically.
        The interpretation of the relations (\resp predicates) of the input structure
        are preserved.
        Hence, it defines a total (non-functional) relation
        from $\Sigma$-structures to~$\Gamma$-structures
        where~$\Gamma=\Sigma\cup\{\colorV\}$ (providing $\colorV\notin\Sigma$).
\end{description}
We say an $\mathsf{L}$-transduction is \defd{deterministic} if it does not use colouring,
it is \defd{non-deterministic} otherwise.
By definition, deterministic $\mathsf{L}$-transductions define functions.

\section{Transducing the laminar-tree}
\label{sec:transduce laminar tree}
In this section,
we present an overlay \CMSO[2]-transduction
that takes as input a laminar set system
and outputs the laminar tree it induces.

\begin{thm}
    \label{thm:transduce laminar tree}
    Let~$\Sigma$ be an extended vocabulary,
    including a unary set predicate name~$\SET$
    and not including the binary relational symbol~$\ancestor$.
    There exists a non-deterministic \CMSO[2]-transduction~$\tau$
    such that,
    for each laminar set system~$(U,\mathcal{S})$
    represented as the $\{\SET\}$-structure~$\mathbb{S}$
    and inducing the laminar tree~$T$ with ${L(T)=U}$,
    and for each $\Sigma$-structure~$\mathbb{A}$
    with~$\mathbb{S}\sqsubseteq\mathbb{A}$,
    $\tau(\mathbb{A})$ is non-empty
    and every output in~$\tau(\mathbb{A})$
    is equal to some $\{\ancestor\}$-structure $\mathbb{T}$ representing~$T$.
\end{thm}

Since the sets from~$\mathcal{S}$
are precisely the sets of leaves of the subtrees of~$T$,
there is an \MSO-transduction
which is the inverse of the above \CMSO[2]-transduction;
that is to say, given an $\{\ancestor\}$-structure~$\mathbb{T}$ representing the laminar tree, it
outputs the original set system~$\mathbb{S}$.
Namely,
\begin{enumerate}[nosep]
    \item An interpretation $\Phi_{\SET}(S)$
    that is true for a set $S$ when there exists an element $a$ such that an element $x$ is in $S$ if and only if $a$ is an ancestor of $x$.
    \item The filtering $\phi(x)$ keeping leaves only.
\end{enumerate}
\medbreak

We prove~\autoref{thm:transduce laminar tree} in \autoref{ssec:transduce laminar tree},
using the key tools developed in~\autoref{ssec:repr},
that allow us to represent each inner node of~$T$ with a pair of leaves from its subtree,
while keeping the number of inner nodes represented by a given leaf bounded.

\subsection{Inner node representatives}
\label{ssec:repr}
In this section, the root of any rooted tree is always an inner node (unless the tree is a unique node).
We fix a rooted tree~$T$,\xspace%
in which every inner node has at least two children (a necessary assumption that is satisfied by laminar trees),
and we let~$V$ denote the set of its nodes ($V=V(T)$)
and~$L\subseteq V$ denote the subset of its leaves ($L=L(T)$).
\smallbreak

Let~$S\subseteq V\setminus L$,
and let~$(\pi,\sigma)$ be a pair of injective mappings from~$S$ to~$L$.
We say that the pair~$(\pi,\sigma)$ \emph{identifies~$S$}
if for each~$s\in S$, $s$ is the least common ancestor of~$\pi(s)$ and~$\sigma(s)$.
For~$s\in S$ and~$x\in V$,
we say that~$x$ is \emph{$s$-requested in~$(\pi,\sigma)$}
if~$x$ lies on the path from~$\pi(s)$ to~$\sigma(s)$
(namely, on either of the paths from~$\pi(s)$ to~$s$ and from~$\sigma(s)$ to~$s$).
We say that~$x$ is \emph{requested in~$(\pi,\sigma)$} if it is $s$-requested for some~$s$.
The pair~$(\pi,\sigma)$ has \emph{unique request}
if every node~$x$ of~$T$ is requested at most once in~$(\pi,\sigma)$,
\ie, there exists at most one~$s\in S$ such that~$x$ is $s$-requested in~$(\pi,\sigma)$. Note that if $(\pi,\sigma)$ has unique request then the paths from $\pi(s)$ to $\sigma(s)$ are pairwise disjoint for all the $s\in S$.
We now state a few basic observations.
\begin{rem}
    \label{rk:identification basic observations}
    Let~$(\pi,\sigma)$ identifying some subset~$S$ of~$V\setminus L$.
    \begin{enumerate}
        \item\label{rk:identification basic observations:commute}
            The reversed pair $(\sigma,\pi)$ also identifies~$S$
            and has unique request whenever~$(\pi,\sigma)$ does.
        \item\label{rk:identification basic observations:sub}
            If~$S'\subset S$, then $(\pi_{|S'},\sigma_{|S'})$ identifies~$S'$,
            and has unique request if~$(\pi, \sigma)$ does.
        \item\label{rk:identification basic observations:self-request}
            For each~$s\in S$, $s$ is $s$-requested in~$(\pi,\sigma)$.
        \item\label{rk:identification basic observations:disjoint}
            If $(\sigma,\pi)$ has unique request,
            then~$\pi(S)$ and~$\sigma(S)$ are disjoint subsets of~$L$.
    \end{enumerate}
\end{rem}

\begin{lem}
    \label{lem:a reprensants}
    Let~$(\pi,\sigma)$ identifying some subset~$S$ of~$V\setminus L$ with unique request.
    Then for each~$a\in\pi(S)$ 
    the node~$\pi^{-1}(a)$ 
    is the least ancestor of~$a$ 
    which belongs to~$S$.
\end{lem}
\begin{proof}
    Let~$a\in\pi(S)$ and let~$s=\pi^{-1}(a)$.
    By definition, $s$ is an ancestor of~$a$ which belongs to~$S$.
    Let~$y$ be the least ancestor of $a$ that is contained in $S$. As $s$ is an ancestor of $a$ belonging to $S$, $y$ must be a descendant of $s$. Hence, $y$ is $s$-requested in~$(\pi,\sigma)$. Additionally, by \autoref{rk:identification basic observations:self-request} of \autoref{rk:identification basic observations},
    $y$ is $y$-requested in~$(\pi,\sigma)$.
    Since~$(\pi,\sigma)$ has unique request, $y=s$.
    Thus, $s$ is the least ancestor of~$a$ which belongs to~$S$.
\end{proof}

Notice that,
by \autoref{rk:identification basic observations:commute} of \autoref{rk:identification basic observations},
a similar result holds for each~$b\in\sigma(S)$.
It follows that the sets~$\pi(S)$ and~$\sigma(S)$
characterize~$(\pi,\sigma)$.
\begin{lem}
    Let~$S\subseteq V\setminus L$,
    and $(\pi,\sigma)$ and $(\pi',\sigma')$
    be two pairs of injections from~$S$ to~$L$
    identifying~$S$ with unique request.
    If~$\pi(S)=\pi'(S)$ and~$\sigma(S)=\sigma'(S)$
    then $\pi=\pi'$ and~$\sigma=\sigma'$.
\end{lem}
\begin{proof}
    Let~$s\in S$, $a=\pi(s)$, and~$s'=\pi'^{-1}(a)$.
    Both~$s$ and~$s'$ are the least ancestor of~$a$ which belongs to~$S$,
    hence~$s=s'$ and~$\pi'(s)=a$.
    Thus~$\pi=\pi'$.
    By \autoref{rk:identification basic observations:commute} of \autoref{rk:identification basic observations},
    we also obtain that~$\sigma=\sigma'$.
\end{proof}

Let~$A$ and~$B$ be two disjoint subsets of~$L$.
We call such a pair~$(A, B)$ a \emph{bi-colouring}.
We say that~$(A,B)$ \emph{identifies~$S$}
if there exists a pair~$(\pi,\sigma)$ identifying~$S$ with unique request
such that~$\pi(S)=A$ and~$\sigma(S)=B$.
By the previous lemma, for a fixed set $S\subseteq V\setminus L$
the pair $(\pi,\sigma)$ identifying $S$ is unique when it exists.
We will also prove that~$S$ is actually uniquely determined from~$(A,B)$.
Before that, we state the following technical lemma.
\begin{lem}
    \label{lem:node cases in identification}
    Let~$(A,B)$ identify some subset~$S$ of~$V\setminus L$
    through a pair~$(\pi,\sigma)$ of injections having unique request.
    Then,
    for each inner node~$x$,
    exactly one of the three following cases holds:
    \begin{enumerate}
        \item
            $x\notin S$,
            $x$ is not requested in~$(\pi,\sigma)$,
            and for each child~$c$ of~$x$,
            $\card{A\cap V(T_c)}=\card{B\cap V(T_c)}$;
        \item
            $x\notin S$,
            $x$ is requested in~$(\pi,\sigma)$,
            and there exists one leaf~$z\in (A\cup B)\cap V(T_x)$
            such that, for each child~$c$ of~$x$,
            $\card{(A\setminus\{z\})\cap V(T_c)}=\card{(B\setminus\{z\})\cap V(T_c)}$;
        \item
            $x\in S$,
            $x$ is requested in~$(\pi,\sigma)$,
            and there exists two leaves
            $a\in A\cap V(T_x)$ and $b\in B\cap V(T_x)$
            such that, for each child~$c$ of~$x$,
            $\card{(A\setminus\{a\})\cap V(T_c)}=\card{(B\setminus\{b\}\cap V(T_c))}$
            and~$\{a,b\}\nsubseteq V(T_c)$.
    \end{enumerate}
\end{lem}
\begin{proof}
    Let~$x$ be an inner node.
    We consider the set~$S'=S\setminus(V(T_x)\setminus\{x\})$ of all nodes which are not proper descendants of $x$
    and the restrictions~$\pi'$ and~$\sigma'$ of, respectively,~$\pi$ and~$\sigma$ to~$S'$.
    By \autoref{rk:identification basic observations:sub} of \autoref{rk:identification basic observations}
    $(\pi',\sigma')$ identify~$S'$ with unique request.
    We denote~$A'=\pi'(S')$ and~$B'=\sigma'(S')$, thus~$(A',B')$ identify~$S'$.
    Let $A'_x=A'\cap V(T_x)$ and $B'_x=B'\cap V(T_x)$ be the sets of representatives contained in $T_x$ of nodes in $S'$.
    Let~$a$ be an element of~$A'_x$.
    The element~$s_a=\pi^{-1}(a)$ is an ancestor of~$x$
    because it is an ancestor of~$a\in V(T_x)$
    and belongs to~$S'$ whence not to~$V(T_x)\setminus\{x\}$.
    Therefore,~$x$ is~$s_a$-requested.
    As~$(\pi',\sigma')$ has unique request,
    there exists at most one element in~$A'_x$.
    Similarly, $B'_x$ has size at most~$1$.
    Moreover, $A'_x\cap B'_x=\emptyset$
    by~\autoref{rk:identification basic observations:disjoint} of \autoref{rk:identification basic observations}.
    We thus we have three cases:
    \begin{description}
        \item[Case~$A'_x\cup B'_x=\emptyset$]
            Then there is no~$s\in S'$ such that~$\pi(s)\in V(T_x)$ or~$\sigma(s)\in V(T_x)$.
            Hence, $x$ is not requested in~$(\pi,\sigma)$, in particular, $x\notin S$.

            Let~$c$ be a child of~$x$.
            If~$\card{A\cap V(T_{c})}\neq\card{B\cap V(T_c)}$,
            then there exists~$a\in(A\cup B)\cap V(T_{c})$
            such that,
            assuming without loss of generality that~$a\in A$ and denoting~$s_a=\pi^{-1}(a)$,
            $\sigma(s_a)\notin V(T_{c})$.
            Hence, $s_a$ is a proper ancestor of~$c$,
            thus, equivalently, an ancestor of~$x$,
            implying that~$x$ is $s_a$-requested in~$(\pi,\sigma)$.
            This contradicts the above argument.
            So, $\card{A\cap V(T_c)}=\card{B\cap V(T_c)}$ for each child~$c$ of~$x$.
        \item[Case~$A'_x\cup B'_x=\{a\}$]
            Assume, without loss of generality, that~$a\in A$,
            and denote $s_a=\pi^{-1}(a)$ and~$b=\sigma(s_a)$.
            We have that~$s_a$ is a proper ancestor of~$x$,
            since it is the least common ancestor of~$a\in V(T_x)$ and~$b\notin V(T_x)$.
            Thus, $x$ is $s_a$-requested and~$s_a\neq x$ so~$x\notin S$.

            Let~$c$ be a child of~$x$.
            If~$\card{A\cap V(T_{c})}\neq\card{B\cap V(T_c)}$,
            then it means that there exists $z\in(A\cup B)\cap V(T_{c})$,
            such that
            either~$z\in A$ and $s_z=\pi^{-1}(z)\notin V(T_{c})$,
            or~$z\in B$ and~$s_z=\sigma^{-1}(z)\notin V(T_{c})$.
            In both cases, $s_z$ is a proper ancestor of~$c$,
            whence an ancestor of~$x$.
            Thus, $x$ is $s_z$-requested in~$(\pi,\sigma)$.
            However, $x$ is $s_a$-requested in~$(\pi,\sigma)$ which has unique request,
            hence~$s_z=s_a$.
            If~$z\in B$, it follows that~$z=b$, which contradicts the fact~$b\notin V(T_x)$.
            Hence, $z\in A$ and thus, $z=\pi(s_a)=a$.
            Therefore, $\card{(A\setminus\{a\})\cap V(T_{c})}=\card{B\cap V(T_c)}$
            for each child~$c$ of~$x$.
        \item[Case~$A'_x=\{a\}$ and~$B'_x=\{b\}$]
            Let~$s_a=\pi^{-1}(a)$ and~$s_b=\sigma^{-1}(b)$.
            The node $x$ is $s_a$-requested and $s_b$-requested in $(\pi,\sigma)$,
            so, by the unique request property, $s_a=s_b$.
            Since~$s_a$ is the least common ancestor of~$a$ and~$b$,
            it belongs to~$V(T_x)$
            whence~$s_a=x$ implying~$x\in X$.

            Let~$c$ be a child of~$x$.
            If~$\card{A\cap V(T_c)}\neq\card{B\cap V(T_c)}$,
            then there exists~$z\in (A\cup B)\cap V(T_{c})$
            such that either~$z\in A$ and~$s_z=\pi^{-1}(z)\notin V(T_{c})$,
            or~$z\in B$ and~$s_z=\sigma^{-1}(z)\notin V(T_{c})$.
            In both cases, $s_z$ is a proper ancestor of~$c$,
            whence an ancestor of~$x$,
            and thus $x$ is $s_z$-requested in~$(\pi,\sigma)$.
            However, $x$ is $x$-requested in~$(\pi,\sigma)$ which has unique request,
            hence~$s_z=x$ and thus~$z\in\{a,b\}$.
            Therefore, $\card{(A\setminus\{a\})\cap V(T_{c})}=\card{(B\setminus\{b\})\cap V(T_c)}$
            for each child~$c$ of~$x$.
            Because the least common ancestor of~$a$ and~$b$ is~$x$,
            there is no child~$c$ of~$x$ containing both~$a$ and~$b$ as leaves,
            \ie, $\{a,b\}\nsubseteq V(T_c)$.
    \end{description}
    This concludes the proof of the statement.
\end{proof}

It follows that for each bi-colouring~$(A, B)$,
there exists at most one set~$S$ of inner nodes
identified by~$(A, B)$.
\begin{lem}
    Let~$(A,B)$ be two disjoint subsets of~$L$
    and let~$S$ and~$S'$ be two subsets of~$V\setminus L$.
    If~$(A,B)$ identify both~$S$ and~$S'$, then~$S=S'$.
\end{lem}
\begin{proof}
    We proceed by contradiction and thus assume~$S\neq S'$.
    Let~$s\in S\setminus S'$.
    By \autoref{lem:node cases in identification},
    the number of children~$c$ of~$s$
    for which the set~$(A\cup B)\cap V(T_{c})$ has odd size is~$2$ since~$s\in S$,
    while it should also be less than~$2$ since~$s\notin S'$.
    A contradiction.
\end{proof}

When~$(A,B)$ identifies a set~$S$,
we call \emph{$A$-representative} (\resp \emph{$B$-representative}) of~$s\in S$
the leaf~$\pi(s)\in A$ (\resp $\sigma(s)\in B$),
where~$(\pi,\sigma)$ witnessing that~$(A,B)$ identifies~$S$.
An example of bi-colouring identifying a subset of inner nodes
is given in \autoref{fig:colouringBinaryCotree}.
\begin{figure}
    \includegraphics[width=\textwidth]{colouring.pdf}%
    \caption{%
        Illustration of a bi-colouring~$(A,B)$
        identifying some set $S\subseteq V\setminus L$
        in a binary tree~$T$.
        Leaves from~$A$ are coloured blue,
        leaves from~$B$ are coloured red,
        and inner nodes from~$S$ are marked purple.
        Furthermore, the two paths
        connecting an inner node~$s\in S$
        to its~$A$- and~$B$-representatives
        are coloured blue and red, respectively.%
    }%
    \label{fig:colouringBinaryCotree}
\end{figure}
\bigbreak

Not every set of inner nodes
has a bi-colouring identifying it.
To ensure that such a pair exists,
we consider \emph{thin sets of inner nodes}. While thin sets always have bi-colourings identifying them, it is also guaranteed that the set of inner nodes can be partitioned into only 4 thin sets.
A subset~$X\subseteq V\setminus L$ is \emph{thin}
when, for each~$x\in X$ not being the root,
on the one hand, the parent~$p_x$ of~$x$ does not belong to~$X$,
and on the other hand, $x$ admits at least one sibling (including possible leaves) that does not belong to~$X$.
Having a thin set~$X$ allows to find branches avoiding it.
\begin{lem}
    \label{lem:thin => path avoiding}
    Let~$X\subseteq V\setminus L$ and~$s\in V\setminus X$.
    If~$X$ is thin, then there exists a leaf~$t\in V(T_s)$
    such that the path from~$t$ to~$s$ avoids~$X$
    (\ie, none of the nodes along this path belong to~$X$).
\end{lem}
\begin{proof}
    If~$T_s$ has height~$0$,
    then~$s$ is a leaf
    and taking~$t=s$ trivially gives the expected path.
    Otherwise,
    $s$ is an inner node and,
    because~$X$ is thin,
    $s$ has at least one child~$c_s$ not belonging to~$X$.
    By induction, there is a path from some leaf~$t\in V(T_{c_s})$ to~$c_s$ avoiding~$X$ and,
    since~$s\notin X$,
    this path could be extended into a path from~$t$ to~$s$ avoiding~$X$.
\end{proof}

The previous lemma allows to identify every thin set.
\begin{lem}
    \label{lem:thin => identifiable}
    If~$X$ is a thin set,
    then there exists a pair~$(\pi,\sigma)$
    of injections from~$X$ to~$L$
    that has unique request and that identifies~$X$.
\end{lem}
\begin{proof}
    We proceed by induction on the size of~$X$.
    If~$X=\emptyset$, the result is trivial.
    Let~$n\in\mathbb{N}$ and suppose that for every thin set of size~$n$
    there exists a pair of injections identifying it with unique request.
    Let~$X$ be a thin set of size~$n+1$,
    and let~$s\in X$ be of minimal depth.
    Clearly, $X\setminus\{s\}$ is thin
    and thus there exists, by induction, a pair~$(\pi,\sigma)$
    of injections from~$X\setminus\{s\}$ to~$L$ identifying~$X\setminus\{s\}$ with unique request.
    Since~$X$ is thin and~$s\in X$,
    we can find two distinct children~$c_a$ and~$c_b$ of~$s$
    not belonging to~$X$.%
    \footnote{\label{fn:out-degree at least 2}Remember that every inner node of~$T$ has at least two children (including possible leaves).}
    Then, by \autoref{lem:thin => path avoiding},
    there exists a leaf~$a\in V(T_{c_a})$ (\resp a leaf~$b\in V(T_{c_b})$)
    such that the path from~$a$ to~$c_a$ (\resp from~$b$ to~$c_b$)
    avoids~$X$.
    In particular, for each node~$y$ along these paths,
    since~$y$ has no ancestor that belongs to~$X$ but~$s$,
    $y$ is not requested in~$(\pi,\sigma)$.
    Hence,
    extending~$\pi$ (\resp~$\sigma$) in such a way that,
    besides mapping each~$x\in X\setminus\{s\}$ to~$\pi(x)$ (\resp to~$\sigma(x)$),
    it maps~$s$ to~$a$ (\resp to~$b$),
    we obtain a pair~$(\hat{\pi},\hat{\sigma})$
    of injections from~$X$ to~$L$
    that identifies~$X$ with unique request.
\end{proof}

A family $F=(A_1,B_1),\ldots,(A_n,B_n)$ of bi-colourings
\emph{identifies} a set~$S\subseteq V\setminus L$,
if there exists a partition~$(S_1,\ldots,S_n)$ of~$S$
such that, for each~$i\in[n]$, $(A_i,B_i)$ identifies~$S_i$.
Whenever $S=V\setminus L$ we say that~$F$ \emph{identifies~$T$}.
%
A collection of subsets of~$V\setminus L$ is \emph{thin}
if each of its subsets is thin.
We now show that there exists a thin $4$-partition.
\begin{lem}
    \label{lem:thin 4-partition}
    There exists a thin $4$-partition of~$V\setminus L$.
\end{lem}
\begin{proof}
    We build such a thin $4$-partition as follows.
    First, consider the partition~$(D_e,D_o)$ of~$V\setminus L$
    in which~$D_e$ (\resp $D_o=V\setminus(D_e\cup L)$) is the set of all inner nodes of even (\resp odd) depth.
    Second, arbitrarily fix one child~$c_x$ of~$x$ for each inner node~$x$,
    and consider the set~$C=\{ c_x \mid x\in V\setminus L \}\setminus L$,
    inducing a partition~$(C,\overline{C})$ of~$V\setminus L$ where~$\overline{C}=V\setminus(L\cup C)$.
    By refining these two bi-partitions, we obtain a $4$-partition which is thin by construction.%
    \footref{fn:out-degree at least 2}
\end{proof}

Hence, four bi-colourings are enough to identify~$T$. For an example see \autoref{fig:partitionThinSets}.
\begin{cor}
    \label{cor:4 bi-colouring}
    There exists a family of four bi-colourings identifying~$T$.%
\end{cor}
\begin{proof}
    The result immediately follows from \autoref{lem:thin => identifiable} and \autoref{lem:thin 4-partition}.%
\end{proof}

\begin{figure}
    \includegraphics[width=\textwidth]{partitionThinSets.pdf}%
    \caption{%
        Illustration of a partition of the inner nodes into 4 thin sets indicated by different shapes/fillings.
        The four colourings identifying each thin set are indicated in the table below the leaves.%
    }%
    \label{fig:partitionThinSets}
\end{figure}
\bigbreak

\subsection{The transduction}
\label{ssec:transduce laminar tree}
The goal of this section is to prove \autoref{thm:transduce laminar tree},
that is, to design a \CMSO[2]-transduction
that produces the laminar tree induced by an input laminar set system.
We fix a laminar set system~$(U,\mathcal{S})$,
represented by a $\{\SET\}$-structure~$\mathbb{S}$.
%
Before proving the theorem,
we make the following basic observation.
On $\mathbb{S}$,
we can define two \MSO-formulas~$\descendant(X,Y)$ and~$\child(X, Y)$
expressing that in the laminar tree induced by~$\mathcal{S}$,
$X$ and~$Y$ are nodes and~$X$ is a descendant or a child of~$Y$, respectively:
\begin{align*}
    \descendant(X, Y)&:=
    \SET(X)\land\SET(Y)\land X\subseteq Y;
    \\*
    \child(X, Y)&:=
    \descendant(X, Y)\land X\neq Y\land
    \forall Z\:
    \Big(\big(
    \descendant(Z, Y)
    \land
    Z\neq Y\big)
    \implies
    \descendant(Z, X)
    \Big).
\end{align*}

The key point of our construction consists in defining
a \CMSO[2]-formula
$\repr_{A,B}(a, X)$ which,
assuming a bi-colouring~$(A, B)$ of the universe
modeled as disjoint unary relations
and
identifying a subset~$S$ of inner nodes of~$T$,
is satisfied exactly when~$X\in S$ and~$a$ is its $A$-representative.
\begin{lem}
    \label{lem:repr MSO}
    Let~$(A,B)$ be a bi-colouring of~$L(T)$
    identifying a subset~$S$ of inner nodes of~$T$.
    There exists a \CMSO[2]-formula~$\repr_{A,B}(a, X)$
    that is satisfied exactly
    when~$X$ is an inner node of~$T$
    that belongs to~$S$ and is $A$-represented by~$a$.
\end{lem}
\begin{proof}
    First, we define a formula~$\phi_{A,B}(X)$
    that, under the above assumption,
    is satisfied exactly when~$X$ belongs to~$S$.
    According to~\autoref{lem:node cases in identification},
    this happens \ifof
    $X$ is a node of~$T$
    (\ie,~$\SET(X)$ is satisfied)
    and there exists~$a\in X\cap A$ and~$b\in X\cap B$
    such that for each child~$Z$ of~$X$,
    $\{a,b\}\nsubseteq Z$
    and
    the set~$(Z\setminus\{a,b\})\cap(A\cup B)$ has even size.
    This property is easily expressed in \CMSO[2],
    using the \MSO-formula $\child(X, Y)$ defined previously,
    as well as the predicate $\SET$:
    \begin{align*}
        \phi_{A,B}(X):=\:
        &
        \SET(X)
        \land
        \exists a\exists b\Big[a\in (X\cap A)\land b\in (X\cap B)\land
        \\*&\qquad
        \forall Z\:
        \Big(
        \child(Z, X) \implies \Big(
        \{a,b\}\nsubseteq Z
        \land
        \C[2]\big((Z\setminus\{a,b\})\cap(A\cup B)\big)
        \Big)\Big)\Big].
    \end{align*}

    Now, we can easily define~$\repr_{A,B}(a, X)$
    based on \autoref{lem:a reprensants}:
    \begin{align*}
        \repr_{A,B}(a, X):=\:&
        \phi_{A,B}(X) \land a\in (X\cap A)
        \land
        \forall Z\subsetneq X\:
        \Big(
        a\in Z
        \implies \neg \phi_{A,B}(Z)
        \Big).
    \end{align*}
    This concludes the proof.
\end{proof}

We are now ready to prove the theorem.
\begin{proof}[Proof of \autoref{thm:transduce laminar tree}]
    The \CMSO[2]-transduction is obtained by composing the following atomic \CMSO[2]-transductions.
    The transduction makes use of the formulas~$\repr_{A, B}(a, X)$
    given by \autoref{lem:repr MSO}.
    \begin{enumerate}
        \item Guess a family of four bi-colourings~$(A_i, B_i)_{i\in[4]}$
            identifying~$T$
            (which exists by \autoref{cor:4 bi-colouring}).
        \item Copy the input graph four times,
            thus introducing four binary relations~$(\copyV[i])_{i\in[4]}$
            where $\copyV[i](x, y)$ indicates that~$x$ is the $i$-th copy
            of the original element~$y$.
        \item Filter the universe
            keeping only
            the original elements
            as well as
            the $i$-th copy of each vertex~$a$
            for which there exists~$X$
            such that~$\repr_{A_i, B_i}(a, X)$.
        \item Define the relation~$\ancestor(x, y)$
            so that it is satisfied exactly when
            there exist~$x'$, $X$, $i$, and~$Y$
            such that,
            on the one hand
            $\descendant(Y, X)$
            and
            $\copyV[i](x,x')\land\repr_{A_i, B_i}(x', X)$,
            and,
            on the other hand,
            either~$y$ is an original element and~$Y=\{y\}$,
            or there exists~$y'$ and~$j$ such that
            $\copyV[j](y, y')\land\repr_{A_j, B_j}(y', Y)$.%
    \end{enumerate}
\end{proof}

\section{Transducing modular decompositions}
\label{sec:transduce modular decomposition}
The set of modules of a directed graph is a specific example of a particular type of set system, a ``weakly-partitive set system'' (and a ``partitive set system'' in the case of an undirected graph). In this section, we first give a general \CMSO[2]-transduction to obtain the canonical tree-like decomposition of a weakly-partitive set system from the set system itself. We then show how to obtain the modular decomposition of a graph via a \CMSO[2]-transduction as an application.
\subsection{Transducing weakly-partitive trees}
\label{ssec:transduce partitive tree}
A set system~$(U,\mathcal{S})$ is \defd{weakly-partitive} if
for every two overlapping sets~$X,Y\in\mathcal{S}$,
the sets~$X\cup Y$, $X\cap Y$, $X\setminus Y$, and~$Y\setminus X$ belong to~$\mathcal{S}$.
It is \defd{partitive} if, moreover,
for every two overlapping sets~$X,Y\in\mathcal{S}$,
their symmetric difference, denoted \defd{$X\symdiff Y$},
also belongs to~$\mathcal{S}$.
By extension, a set family~$\mathcal{S}$ is called \defd{weakly-partitive} or \defd{partitive}
whenever $\left(\bigcup\mathcal{S},\mathcal{S}\right)$ is a set system\xspace%
which is weakly-partitive or partitive, respectively
(note these also requires that~$\emptyset\notin\mathcal{S}$, $\bigcup\mathcal{S}\in\mathcal{S}$, and~$\{a\}\in\mathcal{S}$ for every~$a\in\bigcup\mathcal{S}$).

A member of a set system~$(U,\mathcal{S})$ is said to be~\defd{strong}
if it does not overlap any other set from~$\mathcal{S}$.
The sub-family~\defd{$\mathcal{S}_!$} of strong sets of~$\mathcal{S}$ is thus laminar by definition.
Hence, it induces a laminar tree~$T$.
By extension,
we say that~$T$ is \defd{induced by the set system~$(U,\mathcal{S})$} (or simply by~$\mathcal{S}$).
The next result extends \autoref{thm:transduce laminar tree},
by showing that~$T$ can be \CMSO[2]-transduced from~$\mathcal{S}$.
\begin{lem}
    \label{lem:transduce laminar tree induced by weakly-partitive family}
    Let~$\Sigma$ be an extended vocabulary,
    including a set unary predicate name~$\SET$
    and not including the binary relational symbol~$\ancestor$.
    There exists a non-deterministic \CMSO[2]-transduction~$\tau$
    such that,
    for each set system~$(U,\mathcal{S})$
    represented as the $\{\SET\}$-structure~$\mathbb{S}$
    and inducing the laminar tree~$T$ with ${L(T)=U}$,
    and for each $\Sigma$-structure~$\mathbb{A}$
    with~$\mathbb{S}\sqsubseteq\mathbb{A}$,
    $\tau(\mathbb{A})$ is non-empty
    and every output in~$\tau(\mathbb{A})$
    is equal to some $\{\ancestor\}$-structure $\mathbb{T}$ representing $T$.
\end{lem}
\begin{proof}
    On the $\{\SET\}$-structure~$\mathbb{S}$,
    it is routine to define an \MSO-formula~$\phi_{\SSET}(Z)$
    that identifies those subsets~$Z\subseteq U$
    that are strong members of~$\mathcal{S}$:
    \begin{align*}
        \phi_{\SSET}(Z)&:=
        \SET(Z)
        \land
        \forall X\:
        \Big(
        \big( \SET(X) \land X\cap Z\neq\emptyset \big)
        \implies
        \big( X\subseteq Z \lor Z\subseteq X \big)
        \Big).
    \end{align*}
    Hence, we can design an \MSO-interpretation
    that outputs the $\Sigma\cup\{\SSET\}$-structure
    corresponding to~$\mathbb{A}$ equipped with the set unary predicate $\SSET$
    that selects strong members of~$\mathcal{S}$.
    Thus, up to renaming the predicates~$\SSET$ and~$\SET$,
    by \autoref{thm:transduce laminar tree},
    we can produce,
    through a \CMSO[2]-transduction,
    the $\Sigma\cup\{\ancestor\}$-structure $\mathbb{A}\sqcup\mathbb{T}$
    where~$\mathbb{T}$ is the tree-structure modeling the laminar tree~$T$ induced by~$\mathcal{S}_!$
    with~$L(T)=U$
    (once the output obtained, the interpretation drops the set predicate~$\SSET$ which is no longer needed).
\end{proof}

Clearly, the laminar tree~$T$ of a weakly-partitive set system~$(U,\mathcal{S})$ does not characterize $(U,\mathcal{S})$.
However,
as shown by the \hyperref[thm:weakly-partitive tree]{below theorem},
a labeling of its inner nodes and a controlled partial ordering of its nodes
are sufficient to characterize all the sets of~$\mathcal{S}$.
%
For~$Z$ a set equipped with a partial order~$<$
and~$X$ a subset of~$Z$,
we say that~$X$ is a \defd{$<$-interval}
whenever $<$ defines a total order on~$X$
and for every~$a,b\in X$ and every~$c\in Z$,
$a<c<b$ implies~$c\in X$.%
\begin{thmC}[\cite{chein1981partitive,raoThesis}]
    \label{thm:weakly-partitive tree}%
    Let~$\mathcal{S}$ be a weakly-partitive family,
    $\mathcal{S}_!$ be its subfamily of strong sets,
    and~$T$ be the laminar tree it induces.
    There exists a total labeling function~$\lambda$
    from the set~$V(T)\setminus L(T)$ of inner nodes of~$T$
    to the set~$\{\DEGENERATE,\PRIME,\LINEAR\}$,
    and, for each inner node~$t\in\lambda^{-1}(\LINEAR)$,
    a linear ordering~$<_t$ of its children,
    such that
    every inner node having exactly two children is labeled by~$\DEGENERATE$ and the following two conditions are satisfied:
    \begin{itemize}[nosep]
        \item for each~$X\in\mathcal{S}\setminus\mathcal{S}_!$,
            there exists~$t\in V(T)$
            and a subset~$\mathcal{C}$ of children of~$t$
            such that $X=\bigcup_{c\in\mathcal{C}}L(T_c)$
            and
            either~$\lambda(t)=\LINEAR$
            and~$\mathcal{C}$ is a $<_t$-interval,
            or~$\lambda(t)=\DEGENERATE$;
        \item conversely,
            for each inner node~$t$
            and each non-empty subset~$\mathcal{C}$ of children of~$t$,
            if
            either~$\lambda(t)=\LINEAR$
            and~$\mathcal{C}$ is a $<_t$-interval,
            or~$\lambda(t)=\DEGENERATE$,
            then~$\bigcup_{c\in\mathcal{C}}L(T_c)\in\mathcal{S}$.
    \end{itemize}
\end{thmC}
Furthermore, $T$ and $\lambda$ are uniquely determined from~$\mathcal{S}$,
and, for each inner node~$t$ of~$T$ labeled by~$\LINEAR$, only two orders~$<_t$ are possible,
one being the inverse of the other
(indeed, inverting an order~$<$ does preserve the property of being a $<$-interval).
Hence, every weakly-partitive family~$\mathcal{S}$ is characterized
by a labeled and partially-ordered tree~$(T,\lambda,<)$
where~$T$ is the laminar tree induced by the subfamily~$\mathcal{S}_!$ of strong sets of~$\mathcal{S}$,
$\lambda:V(T)\setminus L(T)\to\{\DEGENERATE,\PRIME,\LINEAR\}$ is the labeling function,
and~$<$ is the partial order $\bigcup_{t\in\lambda^{-1}(\LINEAR)}<_t$ over~$V(T)$.
As, up-to inverting some of the $<_t$ orders, $(T,\lambda,<)$ is unique,
we abusively call it \defd{\strong{the} weakly-partitive tree induced by~$\mathcal{S}$}.
Conversely, a weakly-partitive tree characterizes the unique weakly-partitive set system which induced it.

We naturally model a weakly-partitive tree~$(T,\lambda,<)$ of a partitive set system $(U,\mathcal{S})$
by the $\{\ancestor,\DEGENERATE,\betweeness\}$-structure $\mathbb{T}$
of universe $U_{\mathbb{T}}=V(T)$ such that
$\struct{V(T),\ancestor_{\mathbb{T}}}\sqsubseteq\mathbb{T}$ models~$T$ with ${L(T)=U}$,
$\DEGENERATE_{\mathbb{T}}$ is a unary relation
which selects the inner nodes of~$T$ of label~$\DEGENERATE$,
\ie, ${\DEGENERATE_{\mathbb{T}}=\lambda^{-1}(\DEGENERATE)}$,
and~$\betweeness_{\mathbb{T}}$ is a ternary relation selecting triples $(x,y,z)$ satisfying $x<y<z$ or $z<y<x$.
(Although it is possible, through a non-deterministic \MSO-transduction, to define~$<$ from~$\betweeness_{\mathbb{T}}$,
the use of $\betweeness_{\mathbb{T}}$ rather than~$<_{\mathbb{T}}$ ensures unicity of the output weakly-partitive tree.)
The inner nodes of~$T$ that are labeled by~$\LINEAR$ can be recovered through an \MSO-formula as those inner nodes whose children are related by~$\betweeness$.
The inner nodes labeled by~$\PRIME$ can be recovered through an \MSO-formula as those inner nodes that are labeled neither by~$\DEGENERATE$ nor by~$\LINEAR$.
Using \autoref{thm:weakly-partitive tree},
it is routine to design an \MSO-transduction
which takes as input a weakly-partitive tree
and outputs the weakly-partitive set system which induced it.
The inverse \CMSO[2]-transduction is the purpose of the next result. 
\begin{thm}
    \label{thm:transduce weakly-partitive tree}%
    There exists a non-deterministic \CMSO[2]-transduction~$\tau$
    such that,
    for every weakly-partitive set system~$(U,\mathcal{S})$
    represented as the $\{\SET\}$-structure $\mathbb{S}$
    and inducing the weakly-partitive tree $(T,\lambda,<)$
    represented as the $\{\ancestor,\DEGENERATE,\betweeness\}$-structure $\mathbb{T}$,
    we have $\mathbb{T}\in\tau(\mathbb{S})$
    and every output in $\tau(\mathbb{S})$ is a weakly-partitive tree of $(U,\mathcal{S})$.
\end{thm}
\begin{proof}
    By \autoref{lem:transduce laminar tree induced by weakly-partitive family},
    we have a \CMSO[2]-transduction
    which on inputs~$\mathbb{S}$ outputs the $\{\SET,\ancestor\}$-structure $\mathbb{S}\sqcup\mathbb{T'}$,
    where~$\mathbb{T'}$ is the~$\{\ancestor\}$-structure modeling~$T$, with $L(T)=U=U_{\mathbb{S}}$.
    It thus remains to define a second \CMSO[2]-transduction
    which outputs~$\mathbb{T}$ from~$\mathbb{S}\sqcup\mathbb{T'}$.
    Actually, an \MSO-interpretation is enough.
    Indeed, we can define
    an \MSO-formula~$\phi_{\DEGENERATE}$ which selects the inner nodes of~$T$ that are labeled~$\DEGENERATE$ by~$\lambda$,
    and an \MSO-formula~$\phi_{\betweeness}$ which selects triples~$(x,y,z)$ of elements such that~$x<y<z$ or~$z<y<x$.
    In order to define both formulas,
    we use a functional symbol~$\leafset(t)$
    which is interpreted as the set~$L(T_t)$,
    and is clearly \MSO-definable on~$\{\ancestor\}$-structures such as~$\mathbb{T}$.
    Indeed, an element~$z$ belongs to~$\leafset(t)$ \ifof
    it is a leaf and has~$t$ as ancestor.

    According to \autoref{thm:weakly-partitive tree},
    an inner node~$t$ is labeled~$\DEGENERATE$
    \ifof,
    for every two children $s_1,s_2$ of $t$ the set $L(T_{s_1})\cup L(T_{s_2})$ is in $\mathcal{S}$. We can define this in \MSO using the following formula:
        $$\phi_{\DEGENERATE}(x):=\forall y \forall z\Big(\big(\parent(x,y)\land \parent(x,z)\big)\implies \SET\big(\leafset(y)\cup \leafset(z)\big)\Big).$$


    Now, according to \autoref{thm:weakly-partitive tree},
    we have~$x<y<z$ or~$z<y<x$ exactly when~$x$, $y$, and~$z$ are three distinct children of some node~$t$
    and the three following properties are satisfied:
    \begin{enumerate}[nosep]
        \item\label{it:not degenerate} $L(T_x)\cup L(T_z)$ is not a member of~$\mathcal{S}$;
        \item\label{it:not prime}\label{it:x+y-z} there exists a member of~$\mathcal{S}$ which includes both~$L(T_x)$ and~$L(T_y)$ but excludes~$L(T_z)$;
        \item\label{it:betweeness}\label{it:y+z-x} there exists a member of~$\mathcal{S}$ which includes both~$L(T_y)$ and~$L(T_z)$ but excludes~$L(T_x)$.
    \end{enumerate}
    Indeed,
    when these conditions are satisfied, $t$ is necessarily labeled $\LINEAR$
    (\autoref{it:not degenerate} asserts that it is not labeled by $\DEGENERATE$ and, \eg, \autoref{it:not prime} that it is not labeled by $\PRIME$),
    and, among the possible orderings of~$\{x,y,z\}$, only the ones ensuring~$x<y<z$ or~$z<y<x$ are possible
    ($x<z<y$ and~$y<z<x$ are not possible because of \autoref{it:x+y-z},
    and~$y<x<z$ and~$z<x<y$ are not possible because of \autoref{it:y+z-x}).
    Each of these properties is \MSO-definable:
    \begin{align*}
        \phi_{\betweeness}(x,y,z):=&
            \exists t\Big[ \parent(t,x)\land\parent(t,y)\land\parent(t,z)\land 
            \\*
            &\neg\SET\big(\leafset(x)\cup\leafset(z)\big)\land
            \\*
            &\exists S\Big(\SET(S)\land\big(\leafset(x)\cup\leafset(y)\big)\subseteq S\land\leafset(z)\cap S=\emptyset\Big)\land
            \\*
            &\exists S\Big(\SET(S)\land\big(\leafset(z)\cup\leafset(y)\big)\subseteq S\land\leafset(x)\cap S=\emptyset\Big)\Big].
    \end{align*}
    Once~$\phi_{\DEGENERATE}$ and~$\phi_{\betweeness}$ are defined,
    our transduction drops the original predicate~$\SET$
    which is no longer needed in the output~$\mathbb{T}$.
\end{proof}

If~$\mathcal{S}$ is partitive
then the weakly-partitive tree it induces enjoys a simple form,
and is truly unique.
Indeed, the label $\LINEAR$ and, thus, the partial order~$<$, are not needed.%
\begin{thmC}[\cite{chein1981partitive}]
    \label{thm:partitive tree}
    Let~$\mathcal{S}$ be a weakly-partitive family
    and~$(T,\lambda,<)$ be the weakly-partitive tree it induces.
    If~$\mathcal{S}$ is partitive,
    then~$\lambda^{-1}(\LINEAR)=\emptyset$ and~$<$ is empty.
\end{thmC}
Hence, in case of a partitive set systems~$(U,\mathcal{S})$,
we can consider the simpler object~$(T,\lambda)$,
called \defd{the partitive tree induced by~$\mathcal{S}$}
(or \defd{the partitive tree of~$(U,\mathcal{S})$})
in which~$\lambda$ maps $V(T)\setminus L(T)$ to $\{\DEGENERATE,\PRIME\}$.
%
As a direct consequence of \autoref{thm:partitive tree} and of \autoref{thm:transduce weakly-partitive tree},
we can produce, through a \CMSO[2]-transduction,
the partitive tree induced by a partitive set system
and naturally modeled by an $\{\ancestor,\DEGENERATE\}$-structure.
\begin{cor}
    \label{cor:transduce partitive tree}%
    There exists a non-deterministic \CMSO[2]-transduction~$\tau$
    such that,
    for each partitive set system~$(U,\mathcal{S})$
    represented as the $\{\SET\}$-structure~$\mathbb{S}$
    and inducing the partitive tree $(T,\lambda)$
    represented as the $\{\ancestor,\DEGENERATE\}$-structure $\mathbb{T}$,
    we have~$\mathbb{T}\in\tau(\mathbb{S})$
    and every output in $\tau(\mathbb{S})$ is a partitive tree of $(U,\mathcal{S})$.
\end{cor}

\subsection{Application to modular decomposition}\label{sec:modular-dec}
Let~$G$ be a directed graph
and let~$M\subseteq V(G)$.
We say that~$M$ is a \defd{module (of~$G$)}
if for every $u\notin M$ and every $v,w\in M$,
${uv\in E(G)\!\iff\! uw\in E(G)}$
and~${vu\in E(G)\!\iff\! wu\in E(G)}$.
Clearly,
the empty set, $V(G)$, and all the singletons~$\{x\}$ for~$x\in V(G)$
are modules;
they are called \defd{the trivial modules} of~$G$.
We say a non-empty module~$M$ is \defd{maximal}
if it is not properly contained in any non-trivial module.
Let~$M$ and~$M'$ be two disjoint non-empty modules of~$G$.
Considering the edges that go from~$M$ to~$M'$,
namely edges from the set $(M\times M')\cap E(G)$,
we have two possibilities:
either it is empty, or it is equal to~$M\times M'$.
We write~\defd{$\nfullsucc{M}{M'}$} in the former case
and~\defd{$\fullsucc{M}{M'}$} in the latter.
(It is of course possible to have both~$\fullsucc{M}{M'}$ and~$\fullsucc{M'}{M}$.)
%
A~\defd{modular partition} of~$G$
is a partition $\mathcal{P}=\{M_1,\ldots,M_\ell\}$ of~$V(G)$
such that every~$M_i$ is a non-empty module.
A modular partition $\mathcal{P}=\{M_1,\dots,M_\ell\}$ is called \defd{maximal}
if it is non-trivial and every~$M_i$ is strong and maximal.
Note that every graph has exactly one maximal modular partition.

A \defd{modular decomposition} of~$G$
is a rooted tree~$T$
in which the leaves are the vertices of~$G$,
and for each inner node~$t\in T$,
$t$ has at least two children
and the set~$L(T_t)$ is a module of~$G$.
%
In a modular decomposition~$T$ of~$G$,
for each inner node~$t\in V(T)$
with children~$c_1,\ldots,c_r$,
the family $\mathcal{P}_t=\{L(T_{c_1}),\ldots,L(T_{c_r})\}$
is a modular partition of~$G[V(T_t)]$.
When each such partition is maximal,
the decomposition is unique
and it is called \defd{the maximal modular decomposition} of~$G$.
The maximal modular decomposition~$T$ of~$G$ alone
is not sufficient to characterize~$G$.
However, enriching~$T$ with,
for each inner node~$t$ with children~$c_1,\ldots,c_j$,
the information of which pair of modules $\big(L(T_{c_i}),L(T_{c_j})\big)$
is such that $\fullsucc{L(T_{c_i})}{L(T_{c_j}})$,
yields a unique canonical representation of~$G$.
Formally, the \defd{enriched modular decomposition} of~$G$
is the pair~$(T,F)$
where~$T$ is the maximal modular decomposition of~$G$ (with~$L(T)=V(G)$)
and~$F\subset V(T)\times V(T)$
is a binary relation,
that relates a pair~$(s,t)$ of nodes of~$T$,
denoted $st\in F$,
exactly when~$s$ and~$t$ are siblings and~$\fullsucc{L(T_s)}{L(T_t)}$.
The elements of~$F$ are called \defd{$\mathsf{m}$-edges}.

It should be mentioned that the family of all non-empty modules of~$G$
is known to be weakly-partitive (or even partitive when~$G$ is undirected).
In particular, the maximal modular decomposition~$T$ of~$G$
is the laminar tree induced by the family of strong modules.
Hence \autoref{thm:transduce weakly-partitive tree}
could be used to produce a partially-ordered and labeled tree which displays all the modules of~$G$.
However, this weakly-partitive tree is not sufficient
for being able to recover the graph~$G$ from it.
We now prove how to obtain the enriched modular decomposition of~$G$
through a \CMSO[2]-transduction.

To model enriched modular decompositions as relational structures
we use the relational vocabulary $\{\ancestor, \medge\}$
where $\ancestor$ and $\medge$ are two binary relation names.
An enriched modular decomposition $(T,F)$ of a graph $G$
is modeled by the {$\{\ancestor, \medge\}$}-structure $\mathbb{M}$
with universe $U_{\mathbb{M}}=V(T)$,
$\ancestor_{\mathbb{M}}$ being the set of pairs $(s,t)$
for which~$s$ is an ancestor of~$t$ in~$T$,
and $\medge_{\mathbb{M}}$ being the set of all pairs $(s,t)$
such that $st\in F$
(in particular, $s$ and~$t$ are siblings in~$T$).
%
We use \autoref{lem:transduce laminar tree induced by weakly-partitive family}
in order to transduce the maximal modular decomposition of a graph.
\begin{thm}
    \label{thm:transduce modular decomposition}
    There exists a non-deterministic \CMSO[2]-transduction $\tau$
    such that for every directed graph $G$ represented as the $\{\edge\}$-structure $\mathbb{G}$,
    $\tau(\mathbb{G})$ is non-empty
    and every output in~$\tau(\mathbb{G})$ is equal to some $\{\ancestor,\medge\}$-structure $\mathbb{M}$
    representing the enriched modular decomposition~$(T,F)$ of $G$.
\end{thm}
\begin{proof}
    Let $G$ be a graph represented by the $\{\edge\}$-structure $\mathbb{G}$.
    Several objects are associated to~$G$,
    and each of them can be described by a structure:
    \begin{itemize}[nosep]
        \item let~$\mathcal{M}$ be the family of non-empty modules of~$G$
            and let~$\mathbb{S}$ be the $\{\SET\}$-structure
            modeling the weakly-partitive set system $(V(G),\mathcal{M})$
            with~$U_{\mathbb{S}}=V(G)$;
        \item let~$T$ be the laminar tree induced by the weakly-partitive family~$\mathcal{M}$
            and let~$\mathbb{T}$ be the $\{\ancestor\}$-structure
            modeling it with~$U_{\mathbb{T}}=V(T)$ and~$L(T)=V(G)\subset U_{\mathbb{T}}$;
        \item let~$F$ be the $\mathsf{m}$-edge relation,
            namely the subset of~$V(T)\times V(T)$ such that~$(T,F)$
            is the maximal modular decomposition of~$G$,
            and let~$\mathbb{M}$ be the $\{\ancestor,\medge\}$-structure
            modeling it with~$\mathbb{T}\sqsubset\mathbb{M}$ and $U_{\mathbb{T}}=U_{\mathbb{M}}$.
    \end{itemize}
    Our \CMSO[2]-transduction is obtained by composing the three following transductions:
    \begin{itemize}[nosep]
        \item$\tau_1$: an \MSO-interpretation which outputs the $\{\edge,\SET\}$-structure $\mathbb{G}\sqcup\mathbb{S}$ from~$\mathbb{G}$;
        \item$\tau_2$: the non-deterministic \CMSO[2]-transduction given by \autoref{lem:transduce laminar tree induced by weakly-partitive family}
            which produces the $\{\edge,\SET,\ancestor\}$-structure $\mathbb{G}\sqcup\mathbb{S}\sqcup\mathbb{T}$ from $\mathbb{G}\sqcup\mathbb{S}$;
        \item$\tau_3$: an \MSO-interpretation which outputs the $\{\ancestor,\medge\}$-structure $\mathbb{M}$ from $\mathbb{G}\sqcup\mathbb{S}\sqcup\mathbb{T}$.
    \end{itemize}

    In order to define $\tau_1$ it is sufficient to observe that
    there exists an \MSO-formula $\phi_{\SET}(Z)$
    with one monadic free-variable,
    which is satisfied exactly when~$Z$ is a non-empty module of~$G$.
    Then, since~$\tau_2$ is given by \autoref{lem:transduce laminar tree induced by weakly-partitive family},
    it only remains to define~$\tau_3$.
    Given an inner node~$t$,
    we can select, within \MSO,
    the set~$L(T_t)$ of leaves of the subtree rooted in~$t$.
    We thus assume a function~$\leafset$,
    with one first-order free-variable
    which returns the set of leaves of the subtree rooted at the given node.
    Equipped with this function,
    we can define the \MSO-formula $\phi_{\medge}$
    which selects pairs~$(s,r)$ of siblings
    such that~$sr\in F$.
    Remember that this happen exactly when there exist $u\in L(T_s)$ and~$v\in L(T_r)$
    such that~$uv\in E(G)$.
    Hence, $\phi_{\medge}$ could be defined as:
    \begin{align*}
        \phi_{\medge}(s, r):=
            s\neq r
            \:\land\:
            &
            \exists t\:
            \big(\parent(t,s)\land\parent(t,r)\big)\land
            \\*&
            \exists x\exists y\big(x\in\leafset(s)\land y\in\leafset(r)\land \edge(x,y)\big).
    \end{align*}
    Once defined, the \MSO-interpretation~$\tau_3$
    simply drops all non-necessary relations and predicates
    (namely~$\edge$ and~$\SET$)
    and keeps only the~$\ancestor$ and~$\medge$ relations.
\end{proof}

Notice that it is routine to design a deterministic \MSO-transduction which,
given an $\{\ancestor,\medge\}$-structure $\mathbb{M}$
representing an enriched modular decomposition of some directed graph~$G$,
produces the $\{\edge\}$-structure $\mathbb{G}$ representing~$G$.

\paragraph*{Cographs}
Let~$G$ be a graph, $(T,F)$ be its modular decomposition,
and~$(T,\lambda,<)$ be the weakly-partitive tree induced by the (weakly-partitive) family of its modules.
Let~$t$ be an inner node of~$T$, let~$\mathcal{C}$ be its set of children,
and let~$C$ be the graph~$\big(V(T)\setminus L(T), F\big)\left[\mathcal{C}\right]$ induced by~$F$ on the set of children of~$t$.
It can be checked that,
if~$\lambda(t)=\DEGENERATE$ then~$C$ is either a clique or an independant,
and if~$\lambda(t)=\LINEAR$ then~$C$ is a \emph{tournament consistent with~$<_t$}
(\ie, for every $x,y\in\mathcal{C}$, $xy$ is an edge of~$C$ \ifof $x<_t y$)
or with the inverse of~$<_t$.
In the former case,
we can refine the $\DEGENERATE$ label into~$\SERIES$ and~$\PARALLEL$ labels,
thus expressing that~$C$ is a clique or an independent, respectively.
In the latter case, up-to reversing~$<_t$,
we can ensure that~$C$ is a tournament consistent with~$<_t$.
This yields a refined weakly-partitive tree~$(T,\gamma,<)$,
where~$\gamma$ maps inner nodes to $\{\SERIES,\PARALLEL,\PRIME,\LINEAR\}$
and~$<$ is the order $\bigcup_{t\in\gamma^{-1}(\LINEAR)}<_t$
which ensures that tournaments are consistent with the corresponding $<_t$.
Notice that this labeled and partially-ordered tree is now uniquely determined from~$G$.
Moreover, edges from~$F$ that connect children of a node not labeled by~$\PRIME$
can be recovered from the so-refined weakly-partitive tree.
In particular, if no nodes of~$T$ is labeled by~$\PRIME$,
$(T,F)$ and thus~$G$ is fully characterized by~$(T,\gamma,<)$.
Graphs for which this property holds
are known as \defd{directed cographs},
and can be described by the refined weakly-partitive tree, called \defd{cotree}, explained above and formalized in the following statement.
\begin{thm}
    \label{thm:cotree}
    Let~$G$ be a directed cograph and let~$T$ be the laminar tree induced by the family of its strong modules.
    There exists a unique total labeling~$\lambda$ from the set~$V(T)\setminus L(T)$ of inner nodes of~$T$
    to the set $\{\SERIES,\PARALLEL,\LINEAR\}$ of labels,
    and, for each inner node $t\in\lambda^{-1}(\LINEAR)$,
    a unique linear ordering~$<_t$ of its children,
    such that every inner node having exactly two children is labeled $\SERIES$ or $\PARALLEL$,
    and the following condition is satisfied:
    \begin{itemize}[nosep]
        \item for every two leaves~$x$ and~$y$ of~$T$,
            denoting by~$t$ their least common ancestor,
            $xy$ is an edge of~$G$
            \ifof
            either~$t$ is labeled by~$\LINEAR$ and~$x<_ty$,
            or~$t$ is labeled by~$\SERIES$.
    \end{itemize}
\end{thm}

We naturally model cotrees as $\{\ancestor,\SERIES,\epsord\}$-structures as follow.
A cotree $(T,\gamma,<)$ is modeled by~$\mathbb{C}$
where~$U_{\mathbb{C}}=V(T)$, $\SERIES_{\mathbb{C}}=\gamma^{-1}(\SERIES)$, and $\epsord_{\mathbb{C}}=\{(x,y)\mid x<y\}$.
The nodes that are labeled by $\LINEAR$ could be recovered as those inner nodes whose children are related by $\epsord$,
while the nodes that are labeled by $\PARALLEL$ could be recovered as those inner nodes which are labeled neither by $\SERIES$ nor by $\LINEAR$.
Based on \autoref{thm:cotree} and as a consequence of \autoref{thm:transduce modular decomposition},
we can design a \CMSO[2]-transduction which produces the cotree of a cograph~$G$ from~$G$.
\begin{cor}
    \label{thm:transduce cotree}
    There exists a non-deterministic \CMSO[2]-transduction~$\tau$
    such that, for each cograph~$G$ modeled by the $\{\edge\}$-structure~$\mathbb{G}$,
    $\tau(\mathbb{G})$ is non-empty
    and every output in~$\tau(\mathbb{G})$ is equal to some $\{\ancestor,\SERIES,\epsord\}$-structure~$\mathbb{C}$
    representing the cotree~$(T,\gamma,<)$ of~$G$.
\end{cor}


\section{Transducing split and bi-join decompositions}\label{sec:split}
In this section we consider systems of bipartitions. In graph theory, two known systems of bipartitions are splits and bi-joins. Both are instances of so called weakly-bipartitive systems of bipartitions (or partitive systems of bipartitions in case of undirected graphs). In this section, we first give a \CMSO[2]-transduction of the canonical tree-like decomposition of weakly-bipartitive systems of bipartitions and then derive a \CMSO[2]-transductions for split and bi-join decompositions of a graph.
\subsection{Transducing weakly-bipartitive trees}
A \emph{bipartition system} is a pair $(U,\mathcal{B})$ consisting of a finite set $U$, the \emph{universe}, and a family $\mathcal{B}$ of bipartitions of $U$ such that $\{\emptyset,U\}\notin\mathcal{B}$ and
$\{\{a\},U\setminus \{a\}\}\in \mathcal{B}$ for all $a\in U$.

To model bipartition systems, we use the extended vocabulary~$\{\BIPARTITION\}$
where $\BIPARTITION$ is a set predicate name of arity~$2$.
A bipartition system~$(U,\mathcal{B})$ is thus naturally modeled
as the $\{\BIPARTITION\}$-structure $\mathbb{B}$
with universe~$U_{\mathbb{B}}=U$
and interpretation~$\BIPARTITION_{\mathbb{B}}=\mathcal{B}$.

Two bipartitions $\{X,Y\}$ and $\{X',Y'\}$ of a set $U$ \emph{overlap} if the four sets $X\cap X'$, $X\cap Y'$, $Y\cap X'$ and $Y\cap Y'$ are non-empty.
A bipartition family is said to be \emph{laminar} (aka \emph{overlap-free})
if no two bipartitions in $\mathcal{B}$ overlap.
By extension, we call a family $\mathcal{B}$ of bipartition of a set $U$ \emph{laminar}
whenever $(U,\mathcal{B})$ is a bipartition system which is laminar.
Similarly to laminar set systems, we can associate a tree with each laminar bipartition system. The tree $T$, called the \emph{laminar tree induced by $(U,\mathcal{B})$} (or \emph{laminar tree of $(U,\mathcal{B})$}), associated with a laminar bipartition system $(U,\mathcal{B})$ is unrooted, each element of $U$ corresponds to a leaf of $T$ and for each bipartition $\{X,Y\}$ in $\mathcal{B}$ there is exactly one edge $e$ such that $X=L(T_1)$ and $Y=L(T_2)$ for the two connected components $T_1,T_2$ of $T-e$. We remark that there is a unique such tree and each inner node has degree at least three.
Furthermore, the size of the laminar tree is linearly bounded in the size of the universe $U$.


A bipartition system $(U,\mathcal{B})$ is said to be \emph{weakly-bipartitive} if for every two overlapping bipartitions $\{X,Y\},\{X',Y'\}\in \mathcal{B}$, the biparitions $\{X\cup X',Y\cap Y'\}$, $\{X\cup Y',Y\cap X'\}$, $\{Y\cup X', X\cap Y'\}$ and $\{Y\cup Y', X\cap X'\}$ are in $\mathcal{B}$.
It is \emph{bipartitive}, if additionally for every two overlapping bipartitions $\{X,Y\},\{X',Y'\}\in \mathcal{B}$,
the bipartition $\{X\symdiff X', Y\symdiff Y'\}$ is also in $\mathcal{B}$.
By extension, we call a family $\mathcal{B}$ of bipartitions of a set $U$ \emph{weakly-bipartitive} or \emph{bipartitive} whenever $(U,\mathcal{B})$ is a bipartition system which is weakly-bipartitive or partitive, respectively.

A bipartition $\{X,Y\}$ of a bipartition system~$(U,\mathcal{B})$ is said to be \emph{strong} if $\{X,Y\}$ does not overlap with any other bipartition in $\mathcal{B}$. We denote the subfamily of strong bipartitions in $\mathcal{B}$ by $\mathcal{B}_!$ and remark that $\mathcal{B}_!$ is a laminar family.
Hence, it induces a laminar tree~$T$.
By extension,
we say that~$T$ is \defd{induced by the bipartition system~$(U,\mathcal{B})$} (or simply by~$\mathcal{B}$).
As the laminar tree $T$ of a bipartition system~$(U,\mathcal{B})$ is undirected, we model $T$ by the $\{\tedge\}$-structure $\mathbb{T}$ with universe $U_\mathbb{T}=V(T)$ and binary relation $\tedge_\mathbb{T}$ modeling the edge relation of $T$.

The following lemma gives a connection between laminar trees of set systems and laminar trees of bipartition systems. Note that a very similar statement follows from \cite[Lemma 1.14]{raoThesis} and hence we only provide a proof for sake of completeness in the appendix.
\begin{lem}\label{prop:bipartitiveVSpartitive}
    Let $(U,\mathcal{B})$ be a bipartition system and $a\in U$.
    Then $(U, \mathcal{S}_a)$ is a set system where $\mathcal{S}_a:=\{X: a\notin X \text{ and } \{X,U\setminus X\} \in \mathcal{B}\}\cup \{\{a\},U\}$.
    Moreover, the laminar tree of $(U,\mathcal{B})$ can be obtained from the laminar tree of $(U,\mathcal{S}_a)$ by deleting the root, making node $\{a\}$ adjacent to the node $U\setminus\{a\}$ and making the tree undirected.
\end{lem}


\begin{lem}
    \label{lem:transduce laminar tree induced by weakly-bipartitive family}
    Let~$\Sigma$ be an extended vocabulary,
    including a set unary predicate name~$\BIPARTITION$
    and not including the binary relational symbol~$\tedge$.
    There exists a non-deterministic \CMSO[2]-transduction~$\tau$
    such that,
    for each bipartition system~$(U,\mathcal{B})$
    represented as the $\{\BIPARTITION\}$-structure~$\mathbb{B}$
    and inducing the laminar tree~$T$ with ${L(T)=U}$,
    and for each $\Sigma$-structure~$\mathbb{A}$ with~$\mathbb{B}\sqsubseteq\mathbb{A}$,
    $\tau(\mathbb{A})$ is non-empty
    and every output in~$\tau(\mathbb{A})$ is equal to some $\{\tedge\}$-structure $\mathbb{T}$ representing~$T$.
\end{lem}
\begin{proof}
    Let $\mathcal{B}$ be a family of bipartitions of $U$ represented as the extended relational structure~$\mathbb{B}$ and $\mathbb{A}$ a $\Sigma$-structure
    with $\mathbb{B}\sqsubseteq \mathbb{A}$. We first use the following two simple atomic \CMSO[2]-transduction to obtain a set system from $(U,\mathcal{B})$.
    Guess any colouring $A$ for which there is a single $a\in U$ with $A(a)$ being satisfied. Interpret using the following formula to obtain a $\SET$-predicate:
    $$\phi_\SET(X):=\forall a (A(a)\implies (X=\{a\} \lor X=U \lor a\notin X))\land \BIPARTITION(X,U\setminus X).$$
    The resulting set system $(U,\mathcal{S}_a)$ represented by the $\{\SET\}$-structure $\mathbb{S}$ is a set system by \autoref{prop:bipartitiveVSpartitive}. Hence, we can use the transduction from \autoref{lem:transduce laminar tree induced by weakly-partitive family} to obtain the $\Sigma\cup\{\ancestor\}$-structure $\mathcal{A}\sqcup \mathcal{T}$ where $\mathcal{T}$ is the $\{\ancestor\}$-structure representing the laminar tree $T$ of $(U,\mathcal{S}_a)$. We now use the following two atomic transductions to obtain the weakly-bipartitive tree of $(U,\mathcal{B})$ according to \autoref{prop:bipartitiveVSpartitive}. We fist restrict the universe to all nodes which are not the root of $\mathbb{T}$. We then interpret using the following formula to obtain the $\tedge$-predicate which makes the graph undirected and adds an edge from $a$ to the node representing the set $U\setminus \{a\}$ which is the root of $T$ after removing the node corresponding to $U$:
    $$\phi_\tedge(x,y):=\parent(x,y)\lor \parent(y,x)\lor (x=a \land \lnot \exists z \parent(z,y))\lor (y=a \land \lnot \exists z \parent(z,x)).$$
    Finally, the transduction keeps only $\tedge$ and any $Q\in \Sigma$
    as relations for its output.
\end{proof}


Similar to the set system case, for each weakly-bipartitive family $\mathcal{B}$ we can equip the laminar tree with a labelling of its inner nodes and a partial order which
characterizes the family $\mathcal{B}$.


\begin{thmC}[{\cite[Theorem 3]{montgolfierThesis}}]\label{thm:weakly-bipartitive tree}
    Let~$\mathcal{B}$ be a weakly-bipartitive family,
    $\mathcal{B}_!$ be its subfamily of strong bipartitions,
    and~$T$ the laminar tree it induces.
    There exists a total labeling function~$\lambda$ from the set $V(T)\setminus L(T)$
    of the inner nodes of~$T$ to the set $\{\DEGENERATE,\PRIME,\LINEAR\}$, and, for each inner node $t\in \lambda^{-1}(\LINEAR)$, a linear ordering $<_{t}$ of its neighbors, such that every inner node of degree exactly three is labeled by $\DEGENERATE$ and the following conditions are satisfied:
    \begin{itemize}[nosep]
        \item for each bipartition~$\{X,Y\}\in\mathcal{B}\setminus\mathcal{B}_!$,
            there exists~$t\in V(T)\setminus L(T)$
            and a subset~$\mathcal{C}$ of neighbors of~$t$
            such that either $X$ or $Y$ is equal to the set
            $\bigcup_{c\in\mathcal{C}}L(T^t_c)$ and either $\lambda(t)=\LINEAR$ and $\mathcal{C}$ is a $<_t$-interval, or $\lambda(t)=\DEGENERATE$;
        \item conversely,
            for each inner node~$t$
            and each non-empty subset~$\mathcal{C}$, $\mathcal{C}\not= U$ of neighbors of~$t$,
            the bipartition
            $\big\{\bigcup_{c\in\mathcal{C}}L(T^t_c),\bigcup_{c\notin\mathcal{C}}L(T^t_c)\big\}$ is a member of~$\mathcal{B}$ if either $\lambda(t)=\LINEAR$ and $\mathcal{C}$ is $<_t$-interval or $\lambda(t)=\DEGENERATE$.
    \end{itemize}
\end{thmC}
The tree $T$ and its labelling function $\lambda$ are uniquely determined by $\mathcal{B}$. We note that the total orders $<_t$ are uniquely determined by
$\mathcal{B}$ up to inverting and cyclic shifting, \ie, for a linear order $<$ of $X=\{x_1,\dots,x_\ell\}$ with $x_1<\dots<x_\ell$ we call the order $<'$ of $X$
with $x_\ell<'x_1<'\dots<' x_{\ell-1}$ a \defd{cyclic shift of $<$}. To see this, note that in the theorem above the bipartitions obtained from nodes with label
$\LINEAR$ can be equivalently obtained from a $<_t$-interval or the complement of a $<_t$-interval. Indeed, a cyclic shifting of a linear order $<$ maintains the property of being either a $<$-interval or the complement of a $<$-interval.
Neglecting that $<_t$ are only unique up to inverting and cyclic shifting, we represent every weakly-bipartitive family $\mathcal{B}$ by a triple $(T,\lambda,(<_t)_{t\in \lambda^{-1}(\LINEAR)})$ where $T$ is the laminar tree induced by $\mathcal{B}_!$, $\lambda:V(T)\setminus L(T)\to\{\DEGENERATE,\PRIME,\LINEAR\}$ is the labeling function, and $<_t$ is the total order of the neighbours of $t$ described above. Note that in this case $\bigcup_{t\in\lambda^{-1}(\LINEAR)}<_t$ is not a partial order as it is not necessarily transitive.
We call the triple $(T,\lambda,(<_t)_{t\in \lambda^{-1}(\LINEAR)})$ from \autoref{thm:weakly-bipartitive tree} \defd{\strong{the} bipartitive tree} of $\mathcal{B}$.

We model weakly-bipartitive trees~$(T,\lambda,(<_t)_{t\in \lambda^{-1}(\LINEAR)})$ of a weakly-bipartitive family $\mathcal{B}$ of bipartition of $U$
by the $\{\tedge,\DEGENERATE,\betweeness\}$-structure $\mathbb{T}$
of universe $U_{\mathbb{T}}=V(T)$ such that
$\struct{V(T),\tedge_{\mathbb{T}}}\sqsubseteq\mathbb{T}$ models~$T$ with ${L(T)=U}$,
$\DEGENERATE_{\mathbb{T}}$ is a unary relation
which selects all inner nodes of~$T$ of label~$\DEGENERATE$,
\ie, ${\DEGENERATE_{\mathbb{T}}=\lambda^{-1}(\DEGENERATE)}$,
and~$\betweeness_{\mathbb{T}}$ is a relation of arity four selecting all quadruples $(t,x,y,z)$ for which $y$ is between $x$ and $z$ in $<_t$. Note that in order to define $\betweeness$ we essentially choose one representative order among all orders obtained by inverting and cyclic shifting.
Similarly to the weakly-partitive case, we can recover the inner nodes of~$T$ that are labeled by~$\LINEAR$ through an \MSO-formula as those inner nodes whose neighbors are related by~$\betweeness$.
Then $\PRIME$ nodes are the ones which are labeled neither by~$\DEGENERATE$ nor by~$\LINEAR$ which can be obtained by an $\MSO$-transduction.
Furthermore, \autoref{thm:weakly-bipartitive tree},
gives rise to an \MSO-transduction
which takes as input a weakly-bipartitive tree
and outputs the weakly-bipartitive set system which induced it.


We can now extend the transduction from \autoref{lem:transduce laminar tree induced by weakly-bipartitive family} to obtain a \CMSO[2]-transduction producing the weakly-bipartitive tree of a weakly-bipartitive family.
\begin{thm}
    \label{thm:transduce weakly-bipartitive tree}
    There exists a non-deterministic \CMSO[2]-transduction~$\tau$
    such that,
    for every weakly-bipartitive bipartition system~$(U,\mathcal{B})$
    represented as the $\{\BIPARTITION\}$-structure~$\mathbb{B}$,
    $\tau(\mathbb{B})$ is non-empty and every output in~$\tau(\mathbb{B})$
    is equal to some $\{\tedge,\DEGENERATE,\betweeness\}$-structure~$\mathbb{T}$
    representing the weakly-bipartitive tree~$\left(T,\lambda,(<_t)_{t\in \lambda^{-1}(\LINEAR)}\right)$ of~$(U,\mathcal{B})$.
\end{thm}
\begin{proof}
    \autoref{lem:transduce laminar tree induced by weakly-bipartitive family} already provides a \CMSO[2]-transduction
    which on input $\mathbb{B}$ outputs the $\{\BIPARTITION, \tedge\}$-structure $\mathbb{B}\sqcup \mathbb{T}'$,
    where $\mathbb{T}'$ represents the $\{\tedge\}$-structure modeling the laminar tree $T$ of $(U,\mathcal{B})$.
    To complete the proof of the theorem, we provide an \MSO-transduction which outputs $\mathbb{T}$ given input $\mathbb{B}\sqcup \mathbb{T}'$.
    For the transduction we define an \MSO-formula $\phi_{\DEGENERATE}$ which selects the inner nodes of $T$ which are labeled $\DEGENERATE$,
    and a formula $\phi_{\betweeness}$ which selects quadruples of nodes $(t,x,y,z)$ such that either $x<_t y<_t z$ or $z<_t y<_t x$.
    Recall, that the ordering $<_t$ for each inner node $t$ of the neighbours of $t$ are only unique up to inverting the order and cyclic shift.
    In order to select a representative of $<_t$ which is unique only up to inverting the order, we need non-determinism (details follow).
    We use a functional symbol $\leafset(t,s)$ which is interpreted as the set $L(T_s^t)$.
    The functional predicate $\leafset$ clearly can be defined in \MSO as a element~$z$ belongs to $\leafset(t,s)$
    \ifof~$z$ is a leaf and the path from~$z$ to~$s$ does not contain~$t$.

    Observe that by \autoref{thm:weakly-bipartitive tree} an inner node $t$ is labeled $\DEGENERATE$
    \ifof for every pair of neighbours $s_1,s_2$ of $t$ the bipartition $\big\{L(T^t_{s_1})\cup L(T^t_{s_2}),U\setminus \big(L(T^t_{s_1})\cup L(T^t_{s_2})\big)\big\}$ is in $\mathcal{B}$.
    We can define this in \MSO using the following formula:
    \begin{align*}
        \phi_{\DEGENERATE}&(x):=\forall y \forall z\Big(\big(\tedge(x,y)\land \tedge(x,z)\big)\implies\\
        &\BIPARTITION\big(\leafset(x,y)\cup \leafset(x,z), U\setminus (\leafset(x,y)\cup \leafset(x,z))\big)\Big).
    \end{align*}

    By \autoref{thm:weakly-bipartitive tree}, two neighbors~$x$ and~$y$ of~$t$ are consecutive in~$<_t$
    \ifof, the bipartition $\big\{L(T^t_x)\cup L(T^t_y),U\setminus\big(L(T^t_x)\cup L(T^t_y)\big)\big\}$ is a member of~$\mathcal{B}$
    and at least one of~$x$ or~$y$ is not extremal in~$<_t$.
    Hence, fixing, for each $\LINEAR$-labeled node~$t$,
    the two extremals of~$<_t$ is enough to define~$<_t$ on~$\mathbb{B}\sqcup\mathbb{T'}$ in \MSO.
    To this end, we consider a rooted version of~$T$,
    in which the non-deterministically-guessed root~$r$ is not labeled by~$\LINEAR$,
    and we fix the parent~$p$ of each $\LINEAR$-labeled node~$t$ as the minimum of~$<_t$.
    We also guess a set~$M$ of nodes, such that in~$T$ rooted in~$r$,
    exactly one child~$m$ of each $\LINEAR$-labeled node~$t$ with parent $p$
    belongs to~$M$, and furthermore, the pair $\big\{L(T^t_p)\cup L(T^t_m),U\setminus\big(L(T^t_p)\cup L(T^t_m)\big)\big\}$ is a member of~$\mathcal{B}$.
    Thus, $m$ can be considered as the maximum of~$<_t$.
    Equipped with the relation $\ancestor$ (induced by the rooting of~$T$ in~$r$)
    and the colouring~$M$, we are now able to define $\phi_{\betweeness}$ in \MSO.
    Indeed, for two neighbours $x,y$ of a node $t$ labeled $\LINEAR$, with parent~$p$, and child~$m\in M$, it holds that $x<_t y$ \ifof
    there is a set~$X$ containing both~$L(T_p^t)$ and~$L(T_x^t)$, excluding~$L(T_y^t)$ and $L(T_m^t)$, and satisfying $(X,U\setminus X)\in\mathcal{B}$.
    Hence, we can define in \MSO, $x<_t y$ as follows:
    \begin{align*}
        &x<_t y:=\exists p \exists m \Big[\tedge(t,x)\land\tedge(t,y)\land
            \parent(p,t)\land \tedge(t,m) \land M(m) \land
             \exists X\exists Y
             \\*
            &\Big(\BIPARTITION(X,Y)\land\big(\leafset(t,p)\cup \leafset(t,x)\big)\subseteq X\land \big(\leafset(t,y)\cup \leafset(t,m)\big)\subseteq Y\Big)\Big].
    \end{align*}
    Note that any node with a child in $M$ must necessarily be labeled $\LINEAR$ and hence:
    \begin{align*}
        \phi_{\betweeness}(t,x,y,z)&:= (x<_t y<_t z)\lor(z<_t y<_t x).
    \end{align*}
    Finally, the transduction forgets the predicate $\BIPARTITION$ to produce the output $\mathbb{T}$.
\end{proof}
Indeed, considering a bipartitive family $\mathcal{B}$ makes the use of nodes labeled $\LINEAR$ in the associated tree obsolete, as the following theorem states.
\begin{thmC}[{\cite[Theorem 4]{montgolfierThesis}}]
    Let $\mathcal{B}$ be a weakly-biaprtitive family and let \linebreak $(T,\lambda,(<_t)_{t\in \lambda^{-1}(\LINEAR)})$ be the weakly-bipartitive tree it induces. If $\mathcal{B}$ is bipartitive, then $\lambda^{-1}(\LINEAR)=\emptyset$ and all $<_t$ are empty.
\end{thmC}
As a consequence, we can capture any bipartitive family $\mathcal{B}$ by the simpler structure $(T,\lambda)$, called the \emph{bipartitive tree induced by $\mathcal{B}$} (or \emph{or the bipartitive tree of $\mathcal{B}$}), where $\lambda$ is now a labeling function labeling each inner nodes of $T$ by either $\DEGENERATE$ or $\PRIME$.
Hence, we obtain the following as a corollary from \autoref{thm:transduce weakly-bipartitive tree}, where naturally the bipartitive tree $(T,\lambda)$ of a bipartitive family $\mathcal{B}$ is modeled by a $\{\tedge,\DEGENERATE\}$-structure $\mathbb{B}$.
\begin{cor}
    There exists a non-deterministic \CMSO[2]-transduction $\tau$ such that,
    for each bipartitive set system $(U,\mathcal{B})$ represented as the $\{\BIPARTITION\}$-structure $\mathbb{B}$,
    $\tau(\mathbb{B})$ is non-empty and every output in~$\tau(\mathbb{B})$
    is equal to some $\{\tedge,\DEGENERATE\}$-structure $\mathbb{T}$ representing the bipartitive tree $(T,\DEGENERATE)$ of~$(U,\mathcal{B})$.
\end{cor}


\subsection{Application to split decomposition}
Let $G$ be a directed graph.
A \emph{split} in $G$ is a bipartition $\{X,Y\}$ of $V(G)$ such that there are subsets $X^{\inNeighbours},X^{\outNeighbours}\subseteq X$, $Y^{\inNeighbours},Y^{\outNeighbours}\subseteq Y$ such that for all pairs of vertices $x\in X$ and $y\in Y$ the edge $xy\in E(G)$ \ifof either $x\in X^{\outNeighbours}$ and $y\in Y^{\inNeighbours}$ or $x\in Y^{\outNeighbours}$ and $y\in X^{\inNeighbours}$. Note that if $G$ is connected then $X^{\inNeighbours}\cup X^{\outNeighbours}\not=\emptyset$ and $Y^{\inNeighbours}\cup Y^{\outNeighbours}\not=\emptyset$. We call splits $\{X,Y\}$ for which either $|X|=1$ or $|Y|=1$ \emph{trivial splits}.
In a family $\mathcal{B}$ of splits we use the notion of \emph{strong splits} to coincide with the notion of strong bipartitions.

With a split $\{X,Y\}$ in a connected graph $G$ we associate the following decomposition of $G$ into graphs $G_X$ and $G_Y$. Let $v_X^Y,v_Y^X\notin V(G)$ be two new vertices, called \emph{marker vertices}. Define $G_X$ to be the graph obtained by adding $v_X^Y$ to $G[X]$ and edge $xv_X^Y$ for every $x\in X^{\outNeighbours}$ and edge $v_X^Yx$ for every $x\in X^{\inNeighbours}$. The graph $G_Y$ is defined equivalently. Note that this decomposition captures the graph $G$ as $G$ can be obtained as the union of $G_X[X]$ and $G_Y[Y]$ with additional edges $xy$ whenever either $xv_X^Y\in V(G_X)$ and $v_Y^Xy\in V(G_Y)$ or $xv_Y^X\in V(G_Y)$ and $v_X^Yy\in V(G_X)$.
We can recursively decompose $G$ along splits to obtain a decomposition into graphs $G_1,\dots, G_k$. We call such a decomposition of $G$ into graphs $G_1,\dots, G_k$ a \emph{split decomposition} if every graph $G_i$ is either a singleton consisting of an original (not a marker) vertex or does not have any non-trivial split. We call the graphs $G_i$ \emph{components} of the split decomposition. To be able to reconstruct the graph, we represent a split decomposition of $G$ into graphs $G_1,\dots, G_k$ by an \emph{enriched split decomposition} which is a pair $(H,F)$ consisting of a directed graph $H$ with an additional (symmetric) edge relation $F$. The graph $H$ consisting of the disjoint union $G_1,\dots,G_k$. The edge relation $F$ contains an (undirected) edge between every pair of markers corresponding to one step of the decomposition. For a clear distinction, we call the edges in $H$ \emph{$\mathsf{c}$-edges} and the edges in $F$ \emph{$\mathsf{t}$-edges}.
Note that the splits used in the decomposition are in one-to-one correspondence with the $\mathsf{t}$-edges.

We require the following theorem for our transduction.
\begin{thmC}[\cite{cunningham1982decomposition}]\label{thm:splitsBipartitive}
    If $G$ is connected then the family of splits of $G$ is weakly-bipartitive.
\end{thmC}
Hence, there is a unique split decomposition of $G$, which is called the \emph{canonical split decomposition} of $G$, in which we only decompose along strong splits. Note that the tree we obtain by contracting all edges of the enriched canonical split decomposition $(H,F)$ of $G$ (only leaving the $\textsf{t}$-edges) is isomorphic to the bipartitive tree of the family $\splitSet$.

The following (standard) technical lemma is needed both to show that a graph can be recovered from an enriched split decomposition (by an \MSO transduction) and to define the set of edges in the enriched canonical split decomposition $(H,F)$ when transducing $(H,F)$ from $G$.
\begin{lem}[See \eg~Lemma 2.10 in \cite{AdlerKK17}]\label{lem:edgesInSplitDecomposition}
    Let $G$ be a connected, directed graph and $(H,F)$ an enriched split decomposition of $G$. Then the following hold:
    \begin{itemize}
        \item $uv\in E(G)$ \ifof there exists a directed path from $u$ to $v$ in $(H,F)$ on which $\textsf{c}$-edges and $\textsf{t}$-edges alternate;
        \item $v_{X}^Yv_{X'}^{Y'}$ is a $\textsf{c}$-edge in $(H,F)$ for vertices $v_{X}^Y,v_{X'}^{Y'}$ in the same component of $H$ \ifof there is $uu'\in E(G)$ such that $u\in Y$ and $u'\in Y'$.
    \end{itemize}
\end{lem}

To model enriched split decompositions as relational structures we use the relational vocabulary $\{\tedge,\cedge\}$ where $\tedge$ and $\cedge$ are binary relational symbols. Am enriched split decomposition $(H,F)$ of a graph $G$ is the $\{\tedge,\cedge\}$-relation $\mathbb{H}$ with universe $U_\mathbb{H}=V(H)$, $\tedge_\mathbb{H}$ the set of $\textsf{t}$-edges in $(H,F)$ and $\cedge_{\mathbb{H}}$ the set of $\textsf{c}$-edges of $(H,F)$.

We use \autoref{lem:transduce laminar tree induced by weakly-bipartitive family} to transduce the canonical split decomposition of a graph.
\begin{thm}\label{thm:split}
    There exists a non-deterministic \CMSO[2]-transduction $\tau$
    such that for any connected, directed graph $G$ represented as the $\{\edge\}$-structure $\mathbb{G}$;
    $\tau(\mathbb{G})$ is non-empty
    and every output in~$\tau(\mathbb{G})$ is equal to some $\{\tedge,\cedge\}$-structure $\mathbb{H}$
    representing the canonical split decomposition of $G$.
\end{thm}
\begin{proof}
    Let $G$ be a connected, directed graph represented by the $\{\edge\}$-structure $\mathbb{G}$. We associate the following objects with $G$:
    \begin{itemize}
        \item let $\splitSet$ be the family of splits of $G$ and $\mathbb{B}$ the $\{\BIPARTITION\}$-structure modeling the weakly-bipartitive bipartition system $(V(G),\splitSet)$ with $U_\mathbb{B}=V(G)$;
        \item let $T$ be the laminar tree induced by the weakly-bipartitive family $\splitSet$ and let $\mathbb{T}$ be the $\{\tedge\}$-structure modeling it with $U_\mathbb{T}=V(T)$ and $L(T)=V(G)\subset U_\mathbb{T}$;
        \item let $(H,F)$ be the enriched canonical split decomposition of $G$ represented by the $\{\tedge,\cedge\}$-structure $\mathbb{H}$.
    \end{itemize}
    Our \CMSO[2]-transduction is obtained by composing the following transductions:
    \begin{itemize}
        \item $\tau_1$: an \MSO-interpretation which outputs the $\{\edge,\BIPARTITION\}$-structure $\mathbb{G}\sqcup \mathbb{B}$ from $\mathbb{G}$;
        \item $\tau_2$: the non-deterministic \CMSO[2]-transduction from \autoref{lem:transduce laminar tree induced by weakly-bipartitive family} which outputs the $\{\edge,\BIPARTITION, \tedge\}$-structure $\mathbb{G}\sqcup \mathbb{B}\sqcup \mathbb{T}$ on input $\mathbb{G}\sqcup \mathbb{S}$;
        \item $\tau_3$: an \MSO-transduction which produces the $\{\tedge,\cedge\}$-structure $\mathbb{H}$ from $\mathbb{G}\sqcup \mathbb{B}\sqcup \mathbb{T}$.
    \end{itemize}
    To define $\tau_1$ it is sufficient to observe that there is an \MSO-formula $\phi_{\BIPARTITION}(X,Y)$ with two monadic free-variables $X$ and $Y$ such that $\phi_{\BIPARTITION}$ is satisfied exactly when $\{X,Y\}$ is a split in $G$. As $\tau_2$ is constructed in \autoref{lem:transduce laminar tree induced by weakly-bipartitive family}, we are left with describing how to obtain $\tau_3$.
    We first need to find unique representative for the vertices in $V(H)$. Recall that $V(H)$ consists of two marker vertices $v_X^Y$ and $v_Y^X$ for every strong bipartition $\{X,Y\}\in \splitSet$.
    For ease of notation, we use the name $v^{V(G)\setminus\{u\}}_{\{u\}}$ for the leaf $u\in V(G)$ of $T$. Additionally, we
    set $v_s^t:=v_X^Y$ and $v_t^s:=v_Y^X$ where $\{s,t\}\in E(T)$ is the edge for which $L(T_s)=X$ and $L(T_t)=Y$ where $T_s$ is the component of $T-e$ containing $s$ and $T_t$ the component containing $t$. We use the following atomic transductions to ensure unique representatives for all vertices in $V(H)$. We first guess a colouring $R$ which identifies one non-leaf element of $\mathbb{G}\sqcup \mathbb{B}\sqcup \mathbb{T}$ as root. We use this root element to interpret an auxiliary binary relation $\parent$ in the usual way. Now we copy the structure $\mathbb{G}\sqcup \mathbb{B}\sqcup \mathbb{T}$ once introducing the binary relation $\copyV[1]$ in which $\copyV[1](x,y)$ indicates that $x$ is the copy of $y$. Assume that $\{s,t\}$ is an edge in $T$ corresponding to a strong bipartition in $\splitSet$ and $\parent(s,t)$ is satisfied. We use the original node $t$ to represent the marker $v_t^s$ and we use the copy of $t$ to represent the marker $v_s^t$. As every node has a unique parent in a rooted tree, for each marker the chosen representatives are unique. To express this we use the predicate $\marker(x,y,z)$ which expresses that $x$ is the representative of the marker $v_y^z$ defined as follows:
    $$\marker(x,y,z):= \Big(\parent(z,y)\land x=y\Big)\lor \Big(\parent(y,z)\land \copyV[1](x,z)\Big).$$
    We now filter the universe only keeping all representatives of markers (recall that markers $v_{\{v\}}^{V(G)\setminus\{v\}}$ are the original vertices of $G$). We are left with defining the two edge relations $\tedge$ and $\cedge$. To do so, we use the previously defined functional predicate $\leafset(t,s)$.

    First note that markers $v_t^s$ and $v_{t'}^{s'}$ are in the same component of $H$ \ifof $t=t'$. By \autoref{lem:edgesInSplitDecomposition} $v_t^sv_t^{s'}$ is an $\textsf{c}$-edge in $(H,F)$ \ifof there is an edge $uu'\in E(G)$ with $u\in L(T^t_s)$ and $u'\in L(T^t_{s'})$.
    Consequently, we can express $\cedge$ as follows:
    \begin{align*}
        \phi_{\cedge}(x,x'):=&\exists s \exists s' \exists t \exists y \exists y' \Big( \tedge(t,s)\land \tedge(t,s')\land \marker(x,t,s)\land
        \\*
        &\marker(x',t,s')\land
        y\in \leafset(t,s)\land y'\in \leafset(t,s') \land E(y,y') \Big).
    \end{align*}

    On the other hand, the
    $\tedge$-relation in $\mathbb{H}$ is the relation between pairs of markers and hence can be easily (re)-defined using the $\marker$ predicate by:
    $$\phi_{\tedge}(x,x'):=\exists y \exists z \big(\marker(x,y,z)\land \marker(x',z,y)\big).$$
    Finally, the transduction only keeps the relation $\tedge$ and $\cedge$.
\end{proof}

\subsection{Application to bi-join decomposition}
Bi-joins were introduced in \cite{MontgolfierR05,raoThesis} as a generalization of modules and splits in undirected graphs and we follow their notation. A
\emph{bi-join} in an undirected graph $G$ is a bipartition $\{X,Y\}$ of $V(G)$ such that $\card{X}$, $\card{Y}\geq 1$ and there are subsets $X'\subseteq X$ and
$Y'\subseteq Y$ such that $X'\times Y'\cup (X\setminus X')\times (Y\setminus Y')\subseteq E(G)$, and there are no additional edges between $X$ and $Y$ in $G$.

As for splits, a bi-join $\{X,Y\}$ is considered \emph{trivial} if $\card{X}=1$ or $\card{Y}=1$, and a graph is considered \emph{prime wrt bi-join} if it has only
trivial bi-joins. Let us denote by $\joinSet$ the set of bi-joins of a graph $G$.
A proof of the bipartitiveness of bi-joins can be found in
\cite{raoThesis}.

\begin{thmC}[\cite{raoThesis}]\label{thm:join-bipartitive} For every undirected graph $G$, the system $\joinSet$ is bipartitive. 
\end{thmC}

A corollary of \autoref{thm:join-bipartitive} and \autoref{lem:transduce laminar tree induced by weakly-bipartitive family} we get the existence of a \CMSO[2]-transduction taking as input an
undirected graph $G$ represented as the $\{\edge\}$-structure $\mathbb{G}$ and outputting the laminar tree representing the set of strong bi-joins of
$G$. We would like however, as in the case of splits, to output a relational structure that not only represents the laminar tree induced by $\joinSet$, but also allows to
\CMSO[2]-transduce $\mathbb{G}$. Let's now explain how to construct such a structure from the laminar tree induced by $\joinSet$.
There are two ways to compute such a structure and both are similar to the one for the enhanced canonical split decomposition. Let's first explain how to compute the one
from \cite{montgolfierThesis}.

\paragraph*{Transducing the Skeleton graph \cite{montgolfierThesis}}

If $X$ is a subset of $V(G)$, let $\equiv_X$ be the relation on $X\times X$ where $x\equiv_X y$ if $N(x)\cap V(G)\setminus X = N(y)\cap V(G)\setminus X$. One
easily checks that $\equiv_X$ is an equivalence relation and, if $\{X,Y\}$ is a bi-join, then $\equiv_X$ and $\equiv_Y$ both have at most two equivalence
classes. Let $T$ be the bipartitive tree of $\joinSet$, and for a node $u$, let $u_1$ and $u_2$ be the, at most two, equivalence classes of
$\equiv_{L(T_u)}$. For every node $u$ of $T$, let's denote by $G_u$ the graph whose vertex set is $\{v_1,v_2\mid v$ is a neighbor of $u\}$, and there is an edge
$z_iz_j$ in $G_u$ if there is $x\in z_i$ and $y\in z_j$ such that $xy\in E(G)$. The \emph{skeleton graph} associated with $T$ is the graph obtained by taking
first the disjoint union of the graphs $G_u$, for all nodes $u$ of $T$, and then adding $\textsf{t}$-edges between copies of same equivalence classes. It is
worth mentioning that the skeleton graph is uniquely defined from $T$, and since $T$ is unique, up to isomorphism, one can conclude that the skeleton graph is
unique, up to isomorphism. It is routine to check the following.

\begin{lemC}[\cite{montgolfierThesis}] Let $G$ be a connected graph and let $H$ be its skeleton graph. Then the following hold:
  \begin{enumerate}
  \item two vertices $x$ and $y$ of $G$ are adjacent if and only if there is an alternating path between $x$ and $y$ in $H$;
  \item the bipartitive tree of $\joinSet$ can be obtained from $H$ by contracting the graphs $G_u$, except vertices of $G$, and considering the
    $\mathsf{t}$-edges as its set of edges.
  \end{enumerate}
\end{lemC}


We are now ready to transduce the skeleton graph of a graph, the proof of which is similar to the one for the canonical split decomposition.
\begin{thm}\label{thm:bi-join1}
    There exists a non-deterministic \CMSO[2]-transduction $\tau$
    such that for any connected graph $G$ represented as the $\{\edge\}$-structure $\mathbb{G}$,
    $\tau(\mathbb{G})$ is non-empty
    and every output in~$\tau(\mathbb{G})$
    is equal to some $\{\tedge,\cedge\}$-structure $\mathbb{H}$ representing the skeleton graph of $G$.
\end{thm}
\begin{proof}
    Let $G$ be a connected graph represented by the $\{\edge\}$-structure $\mathbb{G}$, $(V(G), \joinSet)$ the bipartitive set system consisting of all bi-joins
    of $G$ represented by the $\{\BIPARTITION\}$-structure $\mathbb{B}$, $H$ its skeleton graph represented by the $\{\tedge,\cedge\}$-structure $\mathbb{H}$ and $T$ represented by the $\{\tedge,\DEGENERATE\}$-structure, the bipartitive tree of $(V(G),\joinSet)$.

    Our first step is to transduce $\mathbb{B}$ from $\mathbb{G}$. For this it is sufficient to observe that we can easily define an \MSO-formula $\phi_{\BIPARTITION}(X,Y)$ with two monadic free-variables $X$ and $Y$ which defines bi-joins $\{X,Y\}$ of $G$.

    Since $\mathbb{B}$ is bipartitive we can apply \autoref{lem:transduce laminar tree induced by weakly-bipartitive family} to transduce $\mathbb{G}\sqcup \mathbb{B}\sqcup \mathbb{T}$ from $\mathbb{G}\sqcup \mathbb{B}$. In the following we describe how to transduce $\mathbb{H}$ from $\mathbb{G}\sqcup \mathbb{B}\sqcup \mathbb{T}$.

    Contrary to the case of the canonical split decomposition, we cannot deterministically deduce, from $T$, the four markers associated with the bi-join displayed by
    $T$. However, since the set $\{L(T_u)\mid u$ is a node of $T\}$ is a laminar set system, one can use \autoref{cor:4 bi-colouring} to
    non-deterministically associate with each node $u$ of $T$ a leaf as done in the proof of \autoref{thm:transduce laminar tree}. For each node $u$, let's denote
    by $u_1$ this representative. Knowing $u_1$ and $u$, one can make a copy of $u$ and relate it to the equivalence of $\equiv_{L(T_u)}$ not containing
    $u_1$. Now that we can create copies, for each node $u$, representing the two equivalence classes of $\equiv_u$, we leave, to the reader, the description of the
    transduction outputting the desired skeleton graph, as one can use the same copying technique as the one in \autoref{thm:split}.
\end{proof}


\paragraph*{Transducing a decomposition of its Seidel switch equivalence class \cite{raoThesis}}

Another way to represent the bipartitive tree and the graph itself is described in \cite{MontgolfierR05,raoThesis} which uses graphs \emph{characteristic
  graphs} as the ones used to describe a graph with its modular decomposition (see \autoref{sec:modular-dec}). Again, let $T$ be the bipartitive tree of
$\joinSet$. For a node $u$ and $x_1,\ldots, x_p$ leaves in the connected components of $T\setminus u$, not two of them in the same connected component, let
$G[x_1,\ldots, x_p]$ be the \emph{characteristic graph} associated with $u$ and $\{x_1,\ldots, x_p\}$. Unfortunately, while in the case of modular decomposition this graph is
unique, and in the case of split decomposition one can choose it in a canonical way in order to describe the edges of $G$, there may be many characteristic
graphs that can be associated with $u$. However, they are all degenerate if and only if $u$ is.

\begin{lemC}[\cite{MontgolfierR05,raoThesis}]\label{lem:property-characteristic-graph} Let $G$ be a graph and let $T$ be the partitive tree of
  $\joinSet$. For every node $u$ of $T$ and every characteristic graph $G_u$ of $u$, $G_u$ is degenerate if and only if $u$ is degenerate.
\end{lemC}

One can moreover show that they are equivalent with respect to \emph{Seidel switch} \cite{raoThesis}. For a graph $H$ and $W$ a subset of $V(H)$, the
\emph{Seidel switch of $H$ at $W$}, denoted by $\overline{H}^W$, is the graph with $V(H)$ as vertex set and edge set
$E(H[W])\cup E(H\setminus W) \cup \{xy\mid x\in W, y\notin W$ and $xy\notin E(H)\}$ \cite{seidel1991}. One checks that if $H'$ is obtained by a Seidel switch on
$W'$ in $\overline{H}^W$, then $H'=\overline{H}^{W\symdiff W'}$ (see for instance \cite[Lemma 11.1]{ehrenfeuchtHR1999}). Now, two graphs $H$ and $H'$ are \emph{Seidel equivalent} if one is isomorphic to a Seidel switch
of the other. The following tells that all the graphs in a Seidel equivalence class have the same set of bi-joins.

\begin{lemC}[\cite{raoThesis}] Let $H$ be a graph and let $W$ be a subset of $V(H)$. Then, $\{X,Y\}$ is a bi-join in $H$ if and only if $\{X,Y\}$ is a
  bi-join of $\overline{H}^W$.
\end{lemC}

It has been proved that Seidel equivalence classes are in bijection with the set of \emph{two-graphs} \cite{seidel1991}. A \emph{two-graph} is a pair
$(V,\Omega)$ where
 $\Omega$ is a family of subsets of size $3$ on the finite set $V$ such that each subset of size $4$ of $V$ contains an even number of subsets in $\Omega$. For
 every undirected graph $H$, one can associate a two-graph $\Omega(H)$ where $\{x,y,z\}$ belongs to $\Omega(H)$ if $x,y$ and $z$ are all pairwise distinct and
 the number of edges in $H[\{x,y,z\}]$ is odd. Conversely, if $(V,\Omega)$ is a two-graph and $x$ is in $V$, then one can construct a graph $H(\Omega,x)$ with vertex set $V$ and edge
 set $\{yz\mid \{x,y,z\}\in \Omega\}$. The following is proved in \cite{ehrenfeuchtHR1999,seidel1991}.

 \begin{thmC}[{\cite[Theorem 11.7]{ehrenfeuchtHR1999}}]\label{thm:seidel-equiv} Let $V$ be a finite set of size at least $3$. Then the following hold:
   \begin{enumerate}
   \item for every two-graph $(V,\Delta)$, the graph $H(\Delta,x)$ is such that $\Omega(H(\Delta,x))=(V,\Delta)$;
   \item two graphs $G$ and $H$ with vertex set $V$ are Seidel equivalent if and only if $\Omega(G)=\Omega(H)$.
   \end{enumerate}
 \end{thmC}


 For two two-graphs $(V_1,\Omega_1)$ and $(V_2,\Omega_2)$ and $v_1 \in V_1$ and $v_2\in V_2$, their \emph{join wrt $v_1$ and $v_2$} is the two-graph $(V_1\cup V_2\setminus
 \{v_1,v_2\},\Omega)$ where $\Omega = \{\{x,y,z\} \subseteq V_1\cup V_2\setminus
 \{v_1,v_2\}\mid \{x,y,z\}\in \Omega_1 \cup \Omega_2$ or $\{x,y,v_1\}\in \Omega_1, \{z,t,v_2\}\in \Omega_2$ for some $t$, or $\{x,y,v_2\}\in \Omega_2, \{z,t,v_1\}\in \Omega_1$ for some $t\}$.

 A corollary of \autoref{thm:seidel-equiv} is that if $\{X,Y\}$ is a bi-join of a graph $H$ with
 $X\times Y\cap E(H) = X'\times Y' \cup (X\setminus X')\times (Y\setminus Y')$, if we let $H[X]$ (resp. $H[Y]$) as defined for split decomposition with the
 marker $v_X^Y$ (resp. $v_Y^X$) adjacent to $X'$ (resp. $Y'$), $\Omega(H)$ is equal to the join of $\Omega(H[X])$ and of $\Omega(H[Y])$ wrt $v_X^Y$ and
 $v_Y^X$. As observed in \cite{raoThesis}, making $v_X^Y$ (resp. $v_Y^X$) adjacent to $X\setminus X'$ (resp. $Y\setminus Y'$) gives the same two-graph. As a corollary, if we associate with each node $u$ of the bipartitive tree $T$ of $G$ any characteristic graph and join all of them using $\textsf{t}$-edges as in the
split decomposition, one obtains a structure, called \emph{Seidel skeleton}, representing the Seidel equivalence class of $G$, which can by definition transduces $G$. We can summarize then the construction of the desired structure with the following lemma.

\begin{lem}\label{lem:seidel-equiv2} Let $T$ be the bipartive tree of $\joinSet$ and let $S(G,T)$ be the Seidel skeleton of $T$. If $S(G,T)$ is considered as
  a split decomposition and $H$ is the graph it defines, then $G$ is Seidel equivalent to $H$.
\end{lem}

A direct corollary is the desired transduction.

\begin{thm}\label{thm:bi-join2}
    There exists a non-deterministic \CMSO[2]-transduction $\tau$
    such that for any connected graph $G$ represented as the $\{\edge\}$-structure $\mathbb{G}$,
    $\tau(\mathbb{G})$ is non-empty
    and every output in~$\tau(\mathbb{G})$ is equal to some $\{\tedge,\cedge\}$-structure $\mathbb{H}$ representing the Seidel skeleton of $G$.
\end{thm}

Notice that the skeleton graph allows to transduce deterministically $G$ while the structure described here transduces $G$ non-deterministically by guessing the
set $W$ such that Seidel switching it gives $G$.

\section{Conclusion}
We provide transductions for obtaining tree-like graph decompositions such as modular decompositions, cotrees, split decompositions and bi-join decompositions
from a graph using \CMSO. This improves upon results of Courcelle \cite{Courcelle96,Courcelle06} who gave such transductions for ordered graphs. It is worth
mentioning that \autoref{thm:transduce weakly-bipartitive tree} can be also used to \CMSO-transduce canonical decompositions of other structures such as Tutte's
decomposition of matroids or generally 
split-decompositions of submodular functions \cite{cunningham1982decomposition} or modular decompositions of $2$-structures \cite{ehrenfeuchtHR1999} or of hypergraphs
\cite{HabibMMZ22}. As shown by the application given in \cite{Courcelle06} for transducing Whitney's isomorphism class of a graph, a line of research is to more
investigate which structures can be \CMSO-transduced from a graph or a set system by using the transductions from \autoref{thm:summary}. Also, naturally, the
question arises whether counting is necessary or whether \MSO is sufficient to transduce such decompositions.

Furthermore, our results include that \CMSO-transducing rank-decompositions is possible for many classes of graphs.



\begin{thm}\label{thm:cor-cmso-rank-dec} Let $\mathcal{C}$ be a class of graphs for which there is, for every $k$, a \CMSO-transduction $\varphi_{\mathcal{C},k}$, that takes as input a
  graph $G$ in $\mathcal{C}$ and outputs a rank-decomposition of $G$ of width at most $f(k)$, for some computable function $f$. Then, for every positive integer
  $k$ and for every class of graphs
  $\mathcal{D}$ whose prime graphs with respect to split-decomposition belong to $\mathcal{C}$, there is a \CMSO-transduction
  $\psi_{\mathcal{D},k}$ that takes as input a graph $G$ in $\mathcal{D}$ and outputs a rank-decomposition of $G$ of width at most $f(k)$. 
\end{thm}

 A corollary of \autoref{thm:cor-cmso-rank-dec} is that, for any fixed
$k$, there is a \CMSO-transduction that computes a rank-decomposition of small width for any graph belonging to a graph class whose prime graphs have sizes
bounded by $k$ or prime graphs have linear rank-width bounded by $k$. Examples of such graph classes are, for instance,
graphs of rank-width $1$ and many $H$-free graph classes, for some graphs $H$, (see
for instance \cite{BrandstadtDLM05} or \cite{Rao07} for prominent such examples). 

  %% the following bibliography is gererated manually for the sake of brevity
  %% only; please use a separate .bib file in your submission

\bibliographystyle{alphaurl}
\bibliography{bib.bib}

\appendix
\section{Deferred proofs of \autoref{sec:split}}
\begin{proof}[Proof of \autoref{prop:bipartitiveVSpartitive}]
    Let $\mathcal{B}$ be a family of bipartitions of a set $U$ and $a\in U$ any element.
    First observe that $(U,\mathcal{S}_a)$ is a set system
    as $\emptyset\notin \mathcal{S}_a$ (since $\{\emptyset,U\}\notin \mathcal{B}$),
    $U\in \mathcal{S}_a$,
    and $\{b\}\in \mathcal{S}_a$ for every $b\in U$ (since~$\{a\}\in\mathcal{S}_a$ and for~$b\neq a$, $\{\{b\},U\setminus \{b\}\}\in \mathcal{B}$).

    To argue the second assertion, assume that $T$ is the laminar tree of the set system $(U,\mathcal{S}_a)$. First observe that in the set system $(U,\mathcal{S}_a)$ the element $a$ is not contained in any set apart from $U$ and $\{a\}$ and hence the leaf $\{a\}$ is a child of the root $U$ of $T$. Additionally, the set $U\setminus \{a\}$ is in $S_a$ as the bipartition $(\{a\},U\setminus \{a\})$ is in $\mathcal{B}$ by definition and hence $U\setminus \{a\}$ is also a child of the root $U$ of $T$. Let $T'$ be the tree obtained from $T$ by removing node $U$, making all edges undirected and adding an edge from node $\{a\}$ to node $U\setminus \{a\}$. Indeed, by our previous observations, $T'$ is an undirected tree. Furthermore, there is a correspondence between elements of $U$ and the leaves of $T'$ (element $b\in U$ corresponds to leaf $\{b\}$ of $T'$). To check that $T'$ is indeed the laminar tree of $T$, we need to verify that for each bipartition $\{X,Y\}$ in $\mathcal{B}$ there is exactly one edge $e$ of $T'$ such that $X=L(T'_1)$ and $Y=L(T'_2)$ for the two connected components $T'_1,T'_2$ of $T'-e$. We say that the edge \emph{$e$ implements the bipartition $\{X,Y\}$}.
    Assume that $\{X,Y\}$ is any bipartition from $\mathcal{B}$. First consider the case that $\{X,Y\}$ is the bipartition $\{\{a\},U\setminus\{a\}\}$. Then $\{X,Y\}$ is clearly implemented by the edge $\{\{a\},U\setminus\{a\}\}$ of $T'$. Now assume that $\{X,Y\}$ is any other bipartition in $\mathcal{B}$ and assume, without loss of generality, that $a\in Y$. Hence, $X\in \mathcal{S}_a$. Let $X'$ be the parent of $X$ in $T$. Note that $X$ must be a proper subset of $U\setminus \{a\}$ as by our assumption $\{X,Y\}$ is not the partition $\{\{a\},U\setminus \{a\}$ and hence $X'$ cannot be $U$. As we only deleted edge incident to $U$ in our construction of $T'$, $\{X,X'\}$ is an edge in $T'$ which clearly implements the bipartition $\{X,Y\}$. Finally, every bipartition $\{X,Y\}\in \mathcal{B}$ is clearly implemented by only one edge as every inner node of the tree $T'$ has degree at least three. Therefore, $T'$ is the laminar tree of $(U,\mathcal{B})$.
\end{proof}

\end{document}
